\documentclass[12pt]{article}
\usepackage[utf8]{inputenc}
\usepackage{graphicx}
\usepackage{animate}
%导入一些用到的宏包
\usepackage{amsmath,bm,amsfonts,amssymb,enumerate,epsfig,bbm,calc,color,amsthm,ifthen,capt-of,multimedia,bm,multirow}
\usepackage{blindtext}
\usepackage{booktabs}
\usepackage{MnSymbol}
\usepackage{url}
\usepackage{natbib}
\usepackage{changepage}
\usepackage{fancyhdr}
\usepackage[flushleft]{threeparttable}
\newtheorem{theorem}{Theorem}
\usepackage{setspace}
\onehalfspacing

\linespread{1.5}
\textwidth=39pc
\textheight=50pc
\oddsidemargin=0pc
\evensidemargin=0pc
\headsep=0pt
%\headheight=.2cm
\topmargin=0pc
\parindent=1.7pc
\parskip=0pt
\usepackage[flushleft]{threeparttable}


%\newtheorem{theorem}{Theorem}
\DeclareMathOperator*{\argmax}{arg\,max}
\DeclareMathOperator*{\argmin}{arg\,min}
\DeclareMathOperator{\pr}{pr}

\newtheorem{lemma}{Lemma}[]


\title{Supplementary Materials for “Weighted Q-learning for optimal DTRs with nonignorable missing covariates"}
\author{Jian Sun \thanks{School of Data Science,  Fudan University, Shanghai 200433, China. Email: jsun19@fudan.edu.cn.},  Bo Fu \thanks{School of Data Science,  Fudan University, Shanghai 200433, China. Email: fu@fudan.edu.cn.}, and Li Su \thanks{MRC Biostatistics Unit, University of Cambridge, Cambridge, CB2 0SR, UK. Email: li.su@mrc-bsu.cam.ac.uk.}}


\date{}


\providecommand{\keywords}[1]{\textbf{\textit{Keywords  }} #1}

\begin{document}

\maketitle

\section{Web Appendix A: Calibrating the sensitivity parameter in the WQ-SA method}

% \iffalse
% \tb{
% The specification of sensitivity parameters in  nonignorable missingness models (and  in related problems of sensitivity analysis for unmeasured confounding) is  an active  area of research \citep{yin2019simulation, franks2019flexible, Sjolander2022}.  %It generally requires prior knowledge and domain expertise to reason about plausible values of the sensitivity parameters. 
% In this section, we first discuss the sign of the sensitivity parameters in the working model (4) of the main text for the WQ-SA method in the context of the MIMIC-III data. Then we turn to the magnitude of the sensitivity parameters and introduce a simulation-based approach to determining a plausible range of the sensitivity parameters, building upon the methods by \cite{yin2019simulation}.}

% \subsection{The sign of the sensitivity parameter in informative monitoring scenarios}
% \tb{The sign of the sensitivity parameter $\gamma_t'$  in the working model (4) for the WQ-SQ method has a relatively straightforward interpretation: for fixed $\bm{h}_t,a_t$, if $\gamma_t' >0$, then a smaller $y_{pse,t}$ leads to a smaller probability to be missing, and vice versa.  In the context of informative monitoring scenarios in the MIMIC-III data, given the covariates and treatments up to stage $t$ (i.e., $\bm{h}_t,a_t$), patients with worse disease conditions (i.e., larger SOFA scores and smaller values of $y_{pse,t}$ when using the negative SOFA score as the final outcome to be maximized) were more likely to be monitored frequently and  less likely to have missing values of $y_{pse,t}$. Therefore, it is reasonable to assume that the sensitivity parameter $\gamma_t'$ is positive in this scenario.}
% \fi

Although the importance of sensitivity analysis with respect to missing data has been recognized in research and  regulatory agencies' guidelines for analysis and reporting of clinical trials \citep{little2012prevention,european2010guideline,european2020CHMP},
calibrating sensitivity parameters in  nonignorable missingness models (and  in related problems of sensitivity analysis for unmeasured confounding) is still a challenging task and an active  area of research \citep{yin2019simulation, franks2019flexible, Sjolander2022}. We first discuss how to determine the sign of the sensitivity parameter using domain knowledge. Then we describe how to calibrate the magnitude of the sensitivity parameter in the WQ-SA methods. 

\subsection{Determining the sign of the sensitivity parameter}

 In practice, prior knowledge and domain expertise often can  be used to reason about some  characteristics of the missing data mechanism and determine the sign of   sensitivity parameters. In the working model (4) of the main text for the WQ-SA method, the sign of the sensitivity parameter $\gamma_t'$ has a relatively straightforward interpretation: for fixed $\bm{h}_t$ and $a_t$, if $\gamma_t' >0$, then a smaller value of $y_{pse,t}$ leads to a smaller missingness probability, and vice versa. For example, in the MIMIC-III data, given the covariates and treatments up to stage $t$ (i.e., $\bm{h}_t$ and $a_t$), patients with worse disease conditions (i.e., smaller values of $y_{pse,t}$) were more likely to be monitored frequently and  less likely to have missing values of $y_{pse,t}$. Therefore, it is reasonable to assume that the sensitivity parameter $\gamma_t'$ is positive in this scenario.    


\subsection{The simulation-based sensitivity analysis approach by Yin and Shi (2019)}
We now turn to the more challenging task of specifying the plausible magnitude of the sensitivity parameters. The research on specifying the sensitivity parameters in nonignorable missingness models   is relatively sparse in the literature. \cite{yin2019simulation} developed a simulation-based  approach to searching for plausible values of the sensitivity parameters in  logistic models for nonignorable missingness probabilities. This simulation-based method was initially introduced by \cite{gelman1996posterior} for assessing the goodness-of-fit with a Bayesian posterior predictive model. It was further employed for model checking with missing and latent data \citep{gelman2005multiple,copas2000meta,copas2005local,daniels2012bayesian}. 
Using K-nearest-neighbor (KNN) distance metrics, \cite{yin2019simulation} compared the similarity between the \textit{observed} outcome  data with \textit{{replications}} of \textit{{observed outcomes}} simulated from the estimated full data distribution and the estimated  missingness probabilities with fixed values of their sensitivity parameters. If the KNN distance metrics had small values, it indicated that the estimated full data distribution and the missingness probability model (including the specified values of the sensitivity parameters) were \textit{compatible} with the observed outcome data, thus the specified sensitivity parameter values were considered to be plausible.  
Inspired by their methods, we propose a simulation-based approach to determining the plausible range of our sensitivity parameters in the WQ-SA method.


%\tp{I think that if we use $\widehat{Y}_{pse,t}$ instead of $Y_{pse,t}$, then $Y_{pse,t}^{obs}$, $Y_{pse,t}^{\gamma_t',imp}$, $Y_{pse,t}^{\gamma_t',**}$, and $Y_{pse,t}^{\gamma_t',obs}$ may also add the 'widehat' notation. And in Step (1) we mentioned 'equation (5) in the main text', where we use $Y_{pse,t}$. } I've added the following paragraph to clarify this point.  \tp{OK}

\subsection{Distinctions between the nonignorable missing outcome settings in Yin and Shi (2019)  and the setting with nonignorable missing pseudo-outcomes}
Before we proceed to describe the proposed approach, we emphasize some distinctions between the nonignorable missing outcome settings examined by \cite{yin2019simulation} and our DTR setting with nonignorable missing pseudo-outcomes.  In typical nonignorable missing outcome settings in \cite{yin2019simulation}, the observed outcomes are directly available from the raw data. However, due to the potential outcome framework used in the DTR setting, the  pseudo-outcomes before the final stage are conditional expectations of  potential outcomes, thus not directly available from the raw data. Rather, we use estimated Q-functions to obtain estimates of the pseudo-outcomes in the backward-induction procedure. However, due to the nonignorable missing covariates in future stages, the estimates of pseudo-outcomes at the current  stage are partially missing as well and we need to apply missing data methods to address the selection bias from this missingness.  In analogy to the typical non-ignorable outcome settings, we can treat the `observed' pseudo-outcome \textit{estimates} as the \textit{`raw observed outcomes'}  in the DTR setting and  compare their distribution with the distribution of the \textit{replications} of  `\textit{observed}' pseudo-outcomes  in our simulated-based approach. This is reasonable because,  as discussed in Section 2.6 of the main text, the pseudo-outcome estimates converge to the pseudo-outcomes when all the models for the Q functions and the missingness probabilities are correctly specified. When describing  the implementation steps in the next section, we still employ the pseudo-outcome notation and  their models as described in the main text.  
 


\subsection{A simulation-based approach to calibrating the magnitude of the sensitivity parameter}
Let $Y_{pse,t}^{obs}$ and $Y_{pse,t}^{mis}$ denote the observed and unobserved values of $Y_{pse,t}$, respectively, in the sample with complete covariates until stage $t$ (i.e., the patients with $\overline{\bm{R}}_t = \bm{1}_t$). Similarly, let $\widehat{Y}_{pse,t}^{obs}$ denote the observed values of $\widehat{Y}_{pse,t}$ in the sample with $\overline{\bm{R}}_t = \bm{1}_t$.  We propose to  evaluate the plausibility of  each pre-specified value of the sensitivity parameter $\gamma_t'$ by comparing the distribution of the \textit{observed} pseudo-outcome estimates $\widehat{Y}_{pse,t}^{obs}$ with that of $Y_{pse,t}^{obs,\gamma_t'}$, which are the \textit{replications}  of \textit{observed} pseudo-outcomes simulated based on the estimated Q-function at stage $t$ and the estimated working missingness model (4) corresponding to   a specified value of $\gamma_t'$.  


Let $\pr(\cdot)$ and $\pr( \cdot\mid\cdot )$ denote the probability density function and the conditional probability density function, respectively. To simulate data with a specified value of   $\gamma_t'$, we first evaluate the density $\pr(y_{pse,t}\mid \bm{h}_t,a_t,\overline{\bm{R}}_t = \bm{1}_t,R_{pse,t}=0)$, which can be based on  the following Lemma. 
\begin{lemma}
    In the sample with complete covariates until stage $t$ (i.e., $\overline{\bm{R}}_t = \bm{1}_t$), the conditional density of the missing part of $Y_{pse,t}$ (i.e., with $R_{pse,t}=0$) can be expressed by
    \begin{align}
    \frac{\pr(y_{pse,t}\mid \bm{h}_t,a_t,\overline{\bm{R}}_t = \bm{1}_t,R_{pse,t}=0)}{\pr(y_{pse,t}\mid \bm{h}_t,a_t,\overline{\bm{R}}_t = \bm{1}_t,R_{pse,t}=1)} = \frac{\exp(\gamma'_t y_{pse,t})}{E\{\exp(\gamma'_t y_{pse,t}) \mid \bm{h}_t,a_t,\overline{\bm{R}}_t = \bm{1}_t, R_{pse,t}=1\}}. \label{pseudo-obs-mis}
\end{align}
\end{lemma}



In equation~\eqref{pseudo-obs-mis}, $\pr(y_{pse,t}\mid \bm{h}_t,a_t,\overline{\bm{R}}_t = \bm{1}_t,R_{pse,t}=1)$ and $E\{\exp(\gamma'_t y_{pse,t}) \mid \bm{h}_t,a_t,\overline{\bm{R}}_t = \bm{1}_t, R_{pse,t}=1\}$ can be estimated by some nonparametric methods such as kernel regressions based on $\widehat{Y}_{pse,t}^{obs}$. 
Then for any given $\gamma_t'$, we can generate samples from  the estimated $\pr(y_{pse,t}\mid \bm{h}_t,a_t,\overline{\bm{R}}_t = \bm{1}_t,R_{pse,t}=0)$   and impute the missing values in $Y_{pse,t}$ for the patients with complete covariates until stage $t$ (i.e., $\overline{\bm{R}}_t=\bm{1}_t$). We denote these imputed  pseudo-outcomes as $Y_{pse,t}^{\gamma_t',imp}$.
Thus we have complete data ${Y}_{pse,t}^{\gamma_t',*} = (\widehat{Y}_{pse,t}^{obs}, Y_{pse,t}^{\gamma_t',imp})$ for the patients with complete covariates until stage $t$.
%, which can be used to estimate the parameters in the  stage-$t$ Q-function under the future-independent missingness assumption. 
%Under Assumption 4, we have $E[Y_{pse,t}\mid \bm{H}_t,a_t] = E[Y_{pse,t}\mid \bm{H}_t,a_t,\overline{\bm{R}}_t=\bm{1}_t]$. Therefore, we can estimate the stage-$t$ Q-function. 
Under the future-independence missingness assumption (Assumption 4), we have $E(Y_{pse,t}\mid \bm{h}_t,a_t) = E(Y_{pse,t}\mid \bm{h}_t,a_t,\overline{\bm{R}}_t=\bm{1}_t)$, that is, the Q function model applies to the patients with $\overline{\bm{R}}_t=\bm{1}_t$.  Therefore, we  fit   the assumed Q-function model to ${Y}_{pse,t}^{\gamma_t',*}$ and   estimate $E(Y_{pse,t}\mid \bm{h}_t,a_t,\overline{\bm{R}}_t=\bm{1}_t)$.
The density of the residuals of  \iffalse the Q-function \fi this conditional mean  model can also be estimated by nonparametric methods. Combining the estimated \iffalse Q-function \fi conditional mean model and the estimated density of the residuals,  we can obtain the conditional distribution $\pr(y_{pse,t}^{\gamma_t',*}\mid \bm{h}_t,a_t,\overline{\bm{R}}_t=\bm{1}_t)$.  

Next, we take a sample  $Y_{pse,t}^{\gamma_t',**}$ from the estimated $\pr(y_{pse,t}^{\gamma_t',*}\mid \bm{h}_t,a_t, \overline{\bm{R}}_t=\bm{1}_t)$  for the patients with complete covariates until stage $t$ (i.e., $\overline{\bm{R}}_t=\bm{1}_t$).
The missingness indicator $R_{pse,t}^{\gamma_t',*}$ given $Y_{pse,t}^{\gamma_t',**}$ and $\gamma_t'$ would be generated  with the probability $\widehat{\pi}(\bm{h}_t,a_t,y_{pse,t}^{\gamma_t',**};\gamma_t') = \left[1+\exp\left\{\widehat{s}'_{\gamma'_t,t}(\bm{h}_t,a_t) +\gamma'_t y_{pse,t}^{\gamma_t',**}\right\}\right]^{-1}$, where $\widehat{s}'_{\gamma'_t,t}(\bm{h}_t,a_t)$ is estimated by equation (5) in the main text. Drop the values of $Y_{pse,t}^{\gamma_t',**}$ for those with $R_{pse,t}^{\gamma_t',*}=0$ and denote the remaining part of $Y_{pse,t}^{\gamma_t',**}$ as $Y_{pse,t}^{\gamma_t',obs}$, which we refer to as the replications of  observed pseudo-outcomes.
%and drop some values with the probability $1-\widehat{\pi}(\bm{h}_t,a_t,y_{pse,t}^{\gamma_t',**};\gamma_t') = \frac{\exp\{\widehat{s}'_{\gamma'_t,t}(\bm{h}_t,a_t)+\gamma'_t y_{pse,t}^{\gamma_t',**} \}}{1+\exp\{\widehat{s}'_{\gamma'_t,t}(\bm{h}_t,a_t) \}\exp(\gamma'_t y_{pse,t}^{\gamma_t',**})}$ to generate replications of  observed pseudo-outcomes, where $\widehat{s}'_{\gamma'_t,t}(\bm{h}_t,a_t)$ is estimated by equation (5) in the main text. Denote the remaining part of $Y_{pse,t}^{\gamma_t',**}$ as $Y_{pse,t}^{\gamma_t',obs}$.}
Finally, we compare the distribution of $Y_{pse,t}^{\gamma_t',obs}$ with that of $\widehat{Y}_{pse,t}^{obs}$ using the Wilcoxon rank-sum test.  A large  p-value of the Wilcoxon rank-sum test indicates that the replications of  observed pseudo-outcomes are similar to the observed pseudo-outcome estimates $\widehat{Y}_{pse,t}^{obs}$. Therefore, the specified value of $\gamma_t'$ is considered to be plausible.  



For each candidate value of $\gamma_t'$, we formalize the proposed simulation-based approach in the following steps:
\begin{enumerate}
    \item Obtain $\exp\{\widehat{s}'_{\gamma'_t,t}(\bm{h}_t,a_t) \}$ using equation (5) in the main text. 
    %(\tp{Can 'in the main text' be dropped?}) \trd{I think we have to keep it}
    \item Based on the observed pseudo-outcome estimates $\widehat{Y}_{pse,t}^{obs}$, estimate $\pr(y_{pse,t}\mid \bm{h}_t,a_t,\overline{\bm{R}}_t = \bm{1}_t, R_{pse,t}=1)$ and $E\{\exp(\gamma'_t y_{pse,t}) \mid \bm{h}_t,a_t,\overline{\bm{R}}_t = \bm{1}_t, R_{pse,t}=1\}$ nonparametrically. Afterwards, we can estimate $\pr(y_{pse,t}\mid \bm{h}_t,a_t,\overline{\bm{R}}_t = \bm{1}_t,R_{pse,t}=0)$ using equation~\eqref{pseudo-obs-mis}.
    \item Impute $Y_{pse,t}^{\gamma_t',imp}$ for the patients with $(\overline{\bm{R}}_t = \bm{1}_t,R_{pse,t}=0)$  by simulating from the estimated $\pr(y_{pse,t}\mid \bm{h}_t,a_t,\overline{\bm{R}}_t = \bm{1}_t,R_{pse,t}=0)$. Denote $Y_{pse,t}^{\gamma_t',*} = (\widehat{Y}_{pse,t}^{obs}, Y_{pse,t}^{\gamma_t',imp})$.
    \item Based on $Y_{pse,t}^{\gamma_t',*}$, estimate $\pr(y_{pse,t}^{\gamma_t',*}\mid \bm{h}_t,a_t,\overline{\bm{R}}_t=\bm{1}_t)$ with the assumed Q-function model and the nonparametric residual density estimator.
%    \item Generate $Y_{pse,t}^{\gamma_t',**}$ from $\pr(y_{pse,t}^{\gamma_t',*}\mid \bm{h}_t,a_t, \overline{\bm{R}}_t=\bm{1}_t)$ for the patients with $\overline{\bm{R}}_t=\bm{1}_t$. Drop some values of $Y_{pse,t}^{\gamma_t',**}$ with the probability $\frac{\exp\{\widehat{s}'_{\gamma'_t,t}(\bm{h}_t,a_t) \}\exp(\gamma'_t y_{pse,t}^{\gamma_t',**})}{1+\exp\{\widehat{s}'_{\gamma'_t,t}(\bm{h}_t,a_t) \}\exp(\gamma'_t y_{pse,t}^{\gamma_t',**})}$. Denote the remaining $Y_{pse,t}^{\gamma_t',**}$ as $Y_{pse,t}^{\gamma_t',obs}$.
    \item Generate $Y_{pse,t}^{\gamma_t',**}$ from the estimated  $\pr(y_{pse,t}^{\gamma_t',*}\mid \bm{h}_t,a_t, \overline{\bm{R}}_t=\bm{1}_t)$ for the patients with $\overline{\bm{R}}_t=\bm{1}_t$. Generate $R_{pse,t}^{\gamma_t',*}$ from $\widehat{\pi}(\bm{h}_t,a_t,y_{pse,t}^{\gamma_t',**};\gamma_t')$. Let $Y_{pse,t}^{\gamma_t',obs}$ denote  $Y_{pse,t}^{\gamma_t',**}$ with $R_{pse,t}^{\gamma_t',*}=1$.
    \item Compare the distribution of $Y_{pse,t}^{\gamma_t',obs}$ and $\widehat{Y}_{pse,t}^{obs}$ using the Wilcoxon rank-sum test.
    \item Repeat Steps (3)-(6) for many Monte Carlo replicates, report the median of the p-values. Let MCR denote the number of Monte Carlo replicates.
\end{enumerate}



  {\em Remark S1.} In the above procedure, Steps (1)-(5) corresponds to Step (ii) in Section 2.2 of \cite{yin2019simulation} but they have been tailored to our proposed WQ-SA method. It involves the specifications for the sensitivity parameter values, the parametric model for the Q-functions, and the exponential tilting form of the working missingness model. Hence, the similarity between $Y_{pse,t}^{\gamma_t',obs}$ and $\widehat{Y}_{pse,t}^{obs}$ can be used to assess both the plausibility of $\gamma_t'$ and the assumed models in the WQ-SA method.     

 {\em Remark S2.} In Step (6) of the above procedure, we evaluated the similarity between $Y_{pse,t}^{\gamma_t',obs}$ and $\widehat{Y}_{pse,t}^{obs}$ using the Wilcoxon rank-sum test. The Wilcoxon rank-sum test, which is also known as the Mann–Whitney $U$ test, is widely used to determine if 
 a two group sample follows the same distribution. It is a nonparametric test that does not depend upon any particular distributional form (or parameters). As suggested by  Remark 3 in \cite{yin2019simulation}, various methods can be employed to measure the similarity between $Y_{pse,t}^{\gamma_t',obs}$ and $\widehat{Y}_{pse,t}^{obs}$. The Wilcoxon rank-sum test performed well in Simulation 3 in Web Appendix G and in  the MIMIC-III data example,  and it can offer an interpretable measure of the similarity between $Y_{pse,t}^{\gamma_t',obs}$ and $\widehat{Y}_{pse,t}^{obs}$. In contrast,  the KNN distance was computationally expensive, and results relied on tuning parameters in our numerical studies. Therefore, we chose  the Wilcoxon rank-sum test as the similarity measure. 
 
 {\em Remark S3.} To reduce the impact of random errors on the p-values caused by the generation of $Y_{pse,t}^{\gamma_t',imp}$ and $Y_{pse,t}^{\gamma_t',**}$, we repeated Steps (3)-(6) for MCR times. We reported the median of the resulting p-values since they had a skewed distribution.   

% \comment{add a remark about the threshold for median p-values} \tp{OK}

 {\em Remark S4.} In analogy to the conventional significance level $0.05$ for hypothesis testing based on p-values, we considered a threshold of the median p-value  0.05 to determine a plausible range in Simulation 3 and MIMIC-III data analysis. Specifically, a value of $\gamma_t'$  would be included in the plausible range if its corresponding  median p-value was greater than 0.05.

 {\em Remark S5.} The simulation-based approach is for determining a plausible range of the sensitivity parameter, rather than serving as a method for obtaining its point estimate. Following the discussion in \cite{yin2019simulation}, we recommend using a set of sensitivity parameters within a range determined by a tolerance level of p-value medians,  instead of a single value of sensitivity parameter,  in the WQ-SA method.


We evaluated the above approach in Simulation 3 in Web Appendix G and applied it to the MIMIC-III data analysis.%, where $Y_{pse,t}$ and $R_{pse,t}$ are replaced by $\widehat{Y}_{pse,t}$ and $\widehat{R}_{pse,t}$, respectively. } %\trd{I've changed to $\widehat{Y}_{pse,t}$ as above. To change it back, Maybe we can clarify this earlier so people don't confuse by the fact that we have to use $\widehat{Y}_{pse,t}$  to replace ${Y}_{pse,t}$. I will have another look on Tuesday  } \tp{I think that if we use $\widehat{Y}_{pse,t}$ instead of $Y_{pse,t}$, then $Y_{pse,t}^{obs}$, $Y_{pse,t}^{\gamma_t',imp}$, $Y_{pse,t}^{\gamma_t',**}$, and $Y_{pse,t}^{\gamma_t',obs}$ may also add the 'widehat' notation. And in Step (1) we mentioned 'equation (5) in the main text', where we use $Y_{pse,t}$. }






\section{Web Appendix B: Asymptotic properties}

The consistency and asymptotic normality of the estimators in Section 2.6 of the main text  can be established under suitable regularity conditions for estimating equations and an additional assumption that the optimal treatment is unique for all subjects at all stages. Note that the parameters are estimated separately using estimating equations recursively at each stage. Therefore, it is natural to establish the asymptotic properties of the stage-specific parameters recursively as well. For simplicity, we will focus on the two-stage setting, but extensions to the general $T$-stage setting follow directly. 

\begin{theorem}
In the 2-stage scenario, if Assumptions 1-4 hold, the Q-functions at stage 2 are Lipschitz continuous and correctly specified, and $\bm{M}_2(\bm{\theta}_2) = \bm{0}$ has a unique solution, %then
\begin{align*}
    \widehat{\bm{\theta}}_2 - \bm{\theta}_{2}^* &\stackrel{d}{\longrightarrow} N(\bm{0}, \Sigma_{\bm{\theta}_2}) , \\
    \widehat{d}_2^{\mathrm{opt}}(\bm{h}_2) & \stackrel{p}{\longrightarrow} d_2^{\mathrm{opt}}(\bm{h}_2) \text{ for all } \bm{h}_2 \in \mathcal{H}_2,\\
    \widehat{Y}_{pse,1}  &\stackrel{p}{\longrightarrow} Y_{pse,1},
\end{align*}
where $\bm{\theta}_{2}^*$ stands or the true values of $\bm{\theta}_{2}$. $\widehat{\bm{\theta}}_{2}$, $\widehat{d}_{2}^{\mathrm{opt}}(\cdot)$, and $\widehat{Y}_{pse,1}$ denote the estimators of $\bm{\theta}_{2}$, $d_2^{\mathrm{opt}}(\cdot)$, and $Y_{pse,1}$, respectively. Notations $\stackrel{d}{\longrightarrow}$ and  $\stackrel{p}{\longrightarrow}$ stand for converge in distribution and converge in probability, respectively. The detailed form of $\Sigma_{\bm{\theta}_2}$ is given in Web Appendix  C.
\end{theorem}



\begin{theorem}
In the two-stage scenario, if Assumptions 1-5 hold, all models are Lipschitz continuous and correctly specified, and the estimating equations for the Q-functions and the estimating equation \iffalse \eqref{2.2} \fi (3)  in the main text each have a unique solution, then under certain additional regularity conditions, 
\begin{align*}
    \widehat{\bm{\theta}}_{1,EE} - \bm{\theta}_{1}^* &\stackrel{d}{\longrightarrow} N(\bm{0}, \Sigma_{\bm{\theta}_1}) , \\
    \widehat{d}_{1,EE}^{\mathrm{opt}}(\bm{h}_1) & \stackrel{p}{\longrightarrow} d_1^{\mathrm{opt}}(\bm{h}_1) \text{ for all } \bm{h}_1 \in \mathcal{H}_1,
\end{align*}
where $\bm{\theta}_{1}^*$ stands or the true values of $\bm{\theta}_{1}$. $\widehat{\bm{\theta}}_{1,EE}$ and $\widehat{d}_{1,EE}^{\mathrm{opt}}(\cdot)$ denote the estimators of $\bm{\theta}_{1}$ and $d_1^{\mathrm{opt}}(\cdot)$ with the working model \iffalse \eqref{missingnessmodel}\fi (1) in the main text estimated by the EE method, respectively. The regularity conditions and the detailed form of $\Sigma_{\bm{\theta}_1}$ are given in Web Appendix  D.
\end{theorem}


\begin{theorem}
In the two-stage scenario, if Assumptions 1-4 and 6 hold, all models are Lipschitz continuous and correctly specified, the sensitivity parameter $\gamma'$ is correctly specified, and the estimating equations for the Q-functions each have a unique solution, then under certain additional regularity conditions, 
\begin{align*}
    \widehat{\bm{\theta}}_{1,SA} - \bm{\theta}_{1}^* &\stackrel{d}{\longrightarrow} N(\bm{0}, \Sigma'_{\bm{\theta}_1}) , \\
    \widehat{d}_{1,SA}^{\mathrm{opt}}(\bm{h}_1) & \stackrel{p}{\longrightarrow} d_1^{\mathrm{opt}}(\bm{h}_1) \text{ for all } \bm{h}_1 \in \mathcal{H}_1,
\end{align*}
where $\bm{\theta}_{1}^*$ stands or the true values of $\bm{\theta}_{1}$. 
$\widehat{\bm{\theta}}_{1,SA}$ and $\widehat{d}_{1,SA}^{\mathrm{opt}}(\cdot)$ denote the estimators of $\bm{\theta}_{1}$ and $d_1^{\mathrm{opt}}(\cdot)$ with the working model \iffalse \eqref{missingmodel2} \fi (4) in the main text estimated by the SA method, respectively. The regularity conditions and the detailed form of $\Sigma'_{\bm{\theta}_1}$ are given in  Web Appendix  E.
\end{theorem}







\section{Web Appendix C: Proof of Theorem 1}

At stage 2, the parameters $\bm{\theta}_2$ are estimated by solving $\bm{M}_2(\bm{\theta}_2) = \bm{0}$, where $\bm{M}_2(\bm{\theta}_2) =\widehat{E}\left\{\mathbb{I} (\overline{\bm{R}}_{2}=\bm{1}_2) \frac{\partial \left[ \left\{  
  Q_2(\bm{h}_2,a_2;\bm{\theta}_2) - y_{pse,2}\right\} ^2\right]}{\partial \bm{\theta}_2} \right\}$. Under the conditions specified in Theorem 1, $\bm{M}_2(\bm{\theta}_2)$ is an unbiased estimating equation for $\bm{\theta}_2$. With the unbiasedness of these estimating equations, we can get the asymptotic property following \citet[Thm. 6.1]{Newey_McFadden_1994}. Using a first order Taylor expansion centered around $\bm{\theta}_2^*$, we find
\begin{align*}
    \bm{M}_2(\widehat{\bm{\theta}}_2) &= M(\bm{\theta}_2^*) + E\left\{\frac{\partial \bm{M}_2(\bm{\theta}_2^*)}{\partial \bm{\theta}_2}\right\}(\widehat{\bm{\theta}}_2 - \bm{\theta}_2^*) + o(1),\\
    \widehat{\bm{\theta}}_2 - \bm{\theta}_2^* &= -E\left\{\frac{\partial \bm{M}_2(\bm{\theta}_2^*)}{\partial \bm{\theta}_2}\right\}^{-1}\bm{M}_2(\bm{\theta}_2^*)+ o(1).
\end{align*}
For simplicity, we use $\frac{\partial f(\alpha^*, \beta^*)}{\partial \alpha}$ to denote the value of $\frac{\partial f(\alpha,\beta)}{\partial \alpha}$ when $\alpha = \alpha^*, \beta = \beta^*$ for any function $f(\cdot)$ and parameter values $(\alpha^*,\beta^*)$ in the Web Appendix.
\begin{align*}
    E\left[E\left\{\frac{\partial \bm{M}_2(\bm{\theta}_2^*)}{\partial \bm{\theta}_2}\right\}^{-1}\bm{M}_2(\bm{\theta}_2^*)\right] &= \bm{0},\\
    Var\left[E\left\{\frac{\partial \bm{M}_2(\bm{\theta}_2^*)}{\partial \bm{\theta}_2}\right\}^{-1}\bm{M}_2(\bm{\theta}_2^*)\right] &=E\left(\left[E\left\{\frac{\partial \bm{M}_2(\bm{\theta}_2^*)}{\partial \bm{\theta}_2}\right\}^{-1}\bm{M}_2(\bm{\theta}_2^*) \right]^{\otimes 2}\right).
\end{align*}
Besides, $\bm{M}_2(\bm{\theta}_2)$ is defined as an empirical mean of a function of i.i.d. variable. According to the Central Limit Theorem, $(\widehat{\bm{\theta}}_2 - \bm{\theta}_2^*) \stackrel{d}{\longrightarrow} N\left(\bm{0}, E\left\{\left[E\left\{\frac{\partial \bm{M}_2(\bm{\theta}_2^*)}{\partial \bm{\theta}_2}\right\}^{-1}\bm{M}_2(\bm{\theta}_2^*) \right]^{\otimes 2}\right\} \right)$.

The detailed form of $\frac{\partial \bm{M}_2(\bm{\theta}_2^*)}{\partial \bm{\theta}_2}$ can be derived by
\begin{align}
    &\frac{\partial \bm{M}_2(\bm{\theta}_2^*)}{\partial \bm{\theta}_2} = \widehat{E}\left[ \mathbb{I}(\overline{\bm{R}}_{2}=\bm{1}_2) \frac{\partial^2 \left\{Q_2(\bm{h}_2,a_2;\bm{\theta}_2^*) - y_{pse,2}\right\} ^2}{\partial \bm{\theta}_2^2} \right] \notag \\
  &= \widehat{E}\left( 2 \mathbb{I} (\overline{\bm{R}}_{2}=\bm{1}_2)  \frac{\partial \left[ \left\{  
  Q_2(\bm{h}_2,a_2;\bm{\theta}_2^*) - y_{pse,2}\right\} \frac{\partial Q_2(\bm{h}_2,a_2;\bm{\theta}_2^*)}{\partial \bm{\theta}_2} \right]}{\partial \bm{\theta}_2} \right) \notag \\
  &= \widehat{E}\left( 2 \mathbb{I} (\overline{\bm{R}}_{2}=\bm{1}_2) \left[ \left\{ \frac{\partial Q_2(\bm{h}_2,a_2;\bm{\theta}_2^*)}{\partial \bm{\theta}_2} \right\}^2  +  \left\{  
  Q_2(\bm{h}_2,a_2;\bm{\theta}_2^*) - y_{pse,2}\right\} \frac{\partial^2 Q_2(\bm{h}_2,a_2;\bm{\theta}_2^*)}{\partial \bm{\theta}_2^2}   \right] \right), \label{eq1}
\end{align}
where $\frac{\partial Q_2(\bm{h}_2,a_2;\bm{\theta}_2^*)}{\partial \bm{\theta}_2}$ and $ \frac{\partial^2 Q_2(\bm{h}_2,a_2;\bm{\theta}_2^*)}{\partial \bm{\theta}_2^2}$ depend on the specific form of $Q_2(\bm{h}_2,a_2;\bm{\theta}_2)$.



Given that $E(\widehat{\bm{\theta}}_2 - \bm{\theta}_2^*) = 0$, according to the Law of Large Numbers, $\widehat{\bm{\theta}}_2$ converges to $\bm{\theta}_2^*$ in probability. That is, for any $\epsilon > 0$, $P(|\widehat{\bm{\theta}}_2 - \bm{\theta}_2^*| > \epsilon) = 0$ as $n \to \infty$. Furthermore, the Lipschitz continuity of $Q_2(\cdot)$ implies the existence of a Lipschitz constant $L_p > 0$ such that for all ($\bm{\theta}'_2$, $\bm{\theta}_2$), $|Q_2(\bm{h}_2, a_2; \bm{\theta}'_2) - Q_2(\bm{h}_2, a_2; \bm{\theta}_2)| \leq L_p |\bm{\theta}'_2 - \bm{\theta}_2 |$. Therefore, for all $\bm{h}_2 \in \mathcal{H}_2$ and $a_2 \in \{-1, 1\}$, we have
\begin{align*}
    \lim\limits_{n\to\infty}P\{|Q_2(\bm{h}_2,a_2;\widehat{\bm{\theta}}_2) - Q_2(\bm{h}_2,a_2;\bm{\theta}_2^*)|\geq \epsilon\} &\leq  \lim\limits_{n\to\infty} P( L_p |\widehat{\bm{\theta}}_2-\bm{\theta}_2^* | \geq \epsilon)\\
    & =0
\end{align*}
That is, $Q_2(\bm{h}_2,a_2;\widehat{\bm{\theta}}_2)$ converges in probability to $Q_2(\bm{h}_2,a_2;\bm{\theta}_2^*)$ for all fixed $(\bm{h}_2,a_2)$. Since $Q_2(\bm{h}_2,a_2;\bm{\theta}_2)$ can be decomposed as $q_{2,0}(\bm{h}_2; \bm{\beta}_2) + a_2 q_{2,1}(\bm{h}_2;\bm{\psi}_2)$, $q_{2,0}(\bm{h}_2; \widehat{\bm{\beta}}_2)$ and $q_{2,1}(\bm{h}_2;\widehat{\bm{\psi}}_2)$ converge in probability to $q_{2,0}(\bm{h}_2; \bm{\beta}_2^*)$ and $q_{2,1}(\bm{h}_2;\bm{\psi}_2^*)$ for all fixed $\bm{h}_2$, respectively.


Besides, $\widehat{d}_2^{\mathrm{opt}}\left(\bm{h}_2\right) = 2\mathbb{I}\left\{q_{2,1}(\bm{h}_2;\widehat{\bm{\psi}}_2)>0\right\} -1$. Under the standard regularity assumption that $P\left\{q_{2,1}(\bm{h}_2;\widehat{\bm{\psi}}_2)=0\right\} = 0$, we have
\begin{align*}
    P&\left\{\widehat{d}_2^{\mathrm{opt}}(\bm{h}_2) \neq d_2^{\mathrm{opt}}(\bm{h}_2)\right\} = P\left[\mathbb{I}\{q_{2,1}(\bm{h}_2;\widehat{\bm{\psi}}_2)>0\} \neq \mathbb{I}\{q_{2,1}(\bm{h}_2;\bm{\psi}^*_2)>0\}\right]\\
    &= P\{q_{2,1}(\bm{h}_2;\widehat{\bm{\psi}}_2)>0, q_{2,1}(\bm{h}_2;\bm{\psi}_2^*)<0\} +  P\{q_{2,1}(\bm{h}_2;\widehat{\bm{\psi}}_2)<0, q_{2,1}(\bm{h}_2;\bm{\psi}_2^*)>0\}
\end{align*}
Because $\lim\limits_{n\to\infty}q_{2,1}(\bm{h}_t;\widehat{\bm{\psi}}_t) = q_{2,1}(\bm{h}_t;\bm{\psi}_t^*)$, we have $\lim\limits_{n\to\infty}P\left\{q_{2,1}(\bm{h}_2;\widehat{\bm{\psi}}_2)>0, q_{2,1}(\bm{h}_2;\bm{\psi}_2^*)<0\right\} = 0$ and $\lim\limits_{n\to\infty}P\left\{q_{2,1}(\bm{h}_2;\widehat{\bm{\psi}}_2)<0, q_{2,1}(\bm{h}_2;\bm{\psi}_2^*)>0\right\} = 0$. Therefore, for all $\bm{h}_2 \in \mathcal{H}_2$,
$$
\lim\limits_{n\to\infty}P\left\{\widehat{d}_2^{\mathrm{opt}}(\bm{h}_2) \neq d_2^{\mathrm{opt}}(\bm{h}_2)\right\} = 0, \quad 
$$
which means that $\widehat{d}_2^{\mathrm{opt}}(\bm{h}_2) \stackrel{p}{\longrightarrow} d_2^{\mathrm{opt}}(\bm{h}_2)$ for any fixed $\bm{h}_2$. 

Since $Y_{pse,1} = \max _{a_{2} \in \{-1,1\}} Q_{2}\left(\bm{H}_{2}, a_{2};\bm{\theta}_2^*\right)$, we have $y_{pse,1} =  Q_{2}\left\{\bm{h}_{2}, d_2^{\mathrm{opt}}(\bm{h}_2);\bm{\theta}_2^*\right\}$ and $\widehat{y}_{pse,1} =  Q_{2}\left\{\bm{h}_{2}, \widehat{d}_2^{\mathrm{opt}}(\bm{h}_2);\widehat{\bm{\theta}}_2\right\}$ for a given $\bm{h}_2$. Based on the convergence of $\widehat{d}_2^{\mathrm{opt}}(\bm{h}_2)$ and $Q_2(\bm{h}_2,a_2;\widehat{\bm{\theta}}_2)$, we have
\begin{align*}
    \lim\limits_{n\to \infty}& P(|\widehat{y}_{pse,1}-y_{pse,1}|\geq \epsilon) = \lim\limits_{n\to \infty} P\left[\left|Q_2\left\{\bm{h}_2,\widehat{d}_2^{\mathrm{opt}}(\bm{h}_2);\widehat{\bm{\theta}}_2\right\}-Q_2\left\{\bm{h}_2,d_2^{\mathrm{opt}}(\bm{h}_2);\bm{\theta}_2^*\right\}\right|\geq \epsilon \right]\\
    &= \lim\limits_{n\to \infty} P\left[\left|Q_2\left\{\bm{h}_2,\widehat{d}_2^{\mathrm{opt}}(\bm{h}_2);\widehat{\bm{\theta}}_2\right\}-Q_2\left\{\bm{h}_2,d_2^{\mathrm{opt}}(\bm{h}_2);\bm{\theta}_2^*\right\}\right|\geq \epsilon, \widehat{d}_2^{\mathrm{opt}}(\bm{h}_2) \neq d_2^{\mathrm{opt}}(\bm{h}_2) \right] \\
    &\qquad + \lim\limits_{n\to \infty} P\left[ \left|Q_2\left\{\bm{h}_2,\widehat{d}_2^{\mathrm{opt}}(\bm{h}_2);\widehat{\bm{\theta}}_2\right\}-Q_2\left\{\bm{h}_2,d_2^{\mathrm{opt}}(\bm{h}_2);\bm{\theta}_2^*\right\}\right|\geq \epsilon, \widehat{d}_2^{\mathrm{opt}}(\bm{h}_2) = d_2^{\mathrm{opt}}(\bm{h}_2) \right]\\
    &= \lim\limits_{n\to \infty} P\left[ \left|Q_2\left\{\bm{h}_2,d_2^{\mathrm{opt}}(\bm{h}_2);\widehat{\bm{\theta}}_2\right\}-Q_2\left\{\bm{h}_2,d_2^{\mathrm{opt}}(\bm{h}_2);\bm{\theta}_2^*\right\}\right|\geq \epsilon \right]\\
    &=0.
\end{align*}
Therefore, $\widehat{y}_{pse,1} \stackrel{p}{\longrightarrow} y_{pse,1}$ for any fixed $\bm{h}_2$. Subsequently, $\widehat{Y}_{pse,1} \stackrel{p}{\longrightarrow} Y_{pse,1}$.


\section{Web Appendix D: Proof of Theorem 2}

At stage 1, the parameters $\bm{\theta}_1$ are estimated by solving the estimating equation $\bm{M}_1(\bm{\theta}_1) = \bm{0}$, where $\bm{M}_1(\bm{\theta}_1) = \widehat{E} \left(\frac{\widehat{r}_{pse,1}}{\widehat{\pi}_{1}} \mathbb{I}(R_{1}=1) \frac{\partial \left[ \left\{ Q_1(\bm{h}_1,a_1;\bm{\theta}_1) - \widehat{y}_{pse,1}  \right\}^2\right]}{\partial \bm{\theta}_1} \right)$. Following \cite{wallace2015doubly}  and \cite{Simoneau_et_al_2020}, we establish the asymptotic normality of $\widehat{\bm{\theta}}_1$ under the regularity condition that $P(\{q_{2,1}(\bm{H}_2; \widehat{\bm{\psi}}_2) = 0\})=0$. 
Since $(\widehat{Y}_{pse,t},\widehat{R}_{pse,t},\widehat{\pi}_t)$ are calculated with $(\widehat{\bm{\theta}}_2,\widehat{\gamma}_1)$, $\bm{M}_1$ depends $(\widehat{\bm{\theta}}_2,\widehat{\gamma}_1)$ and thus can also be denoted as $\bm{M}_1(\bm{\theta}_1, \widehat{\bm{\theta}}_2,\widehat{\gamma}_1)$. The variance of $\bm{\theta}_1$ must adjust for the plug-in estimates of nuisance parameters $(\bm{\theta}_2,\gamma_1)$. By performing a first-order Taylor expression of $\bm{M}_1(\bm{\theta}_1, \widehat{\bm{\theta}}_2,\widehat{\gamma}_1)$ around $(\bm{\theta}_2^*,\gamma_1^*)$, we have
\begin{align*}
    M_{\mathrm{adj}, 1}(\bm{\theta}_1) = \bm{M}_1(\bm{\theta}_1, \widehat{\bm{\theta}}_2,\widehat{\gamma}_1) + E\left\{ \frac{\partial \bm{M}_1(\bm{\theta}_1, \bm{\theta}_2^*,\gamma_1^*)}{\partial (\bm{\theta}_2,\gamma_1)} \right\}
    \left\{\left(\begin{array}{l}
\widehat{\bm{\theta}}_2 \\
\widehat{\gamma}_1
\end{array}\right)-\left(\begin{array}{l}
\bm{\theta}_2^* \\
\gamma_1^*
\end{array}\right)\right\}+ o(1),
\end{align*}
where $
    E\left\{ \frac{\partial \bm{M}_1(\bm{\theta}_1, \bm{\theta}_2^*,\gamma_1^*)}{\partial (\bm{\theta}_2,\gamma_1)} \right\} = E\left\{ \left( \frac{\partial \bm{M}_1(\bm{\theta}_1, \bm{\theta}_2^*,\gamma_1^*)}{\partial \bm{\theta}_2} , \frac{\partial \bm{M}_1(\bm{\theta}_1, \bm{\theta}_2^*,\gamma_1^*)}{\partial \gamma_1}  \right) \right\}
$, and we assume that $\bm{M}_1$ is infinitely differentiable for $(\bm{\theta}_1,\bm{\theta}_2,\gamma_1)$. Therefore, 
\begin{align*}
    \bm{M}_{adj,1}(\widehat{\bm{\theta}}_1) &= \bm{M}_{adj,1}(\bm{\theta}_1^*) + E\left\{\frac{\partial \bm{M}_{adj,1}(\bm{\theta}_1^*)}{\partial \bm{\theta}_1}\right\}(\widehat{\bm{\theta}}_1 - \bm{\theta}_1^*) + o(1),\\
    \widehat{\bm{\theta}}_1 - \bm{\theta}_1^* &= -E\left\{\frac{\partial \bm{M}_{adj,1}(\bm{\theta}_1^*)}{\partial \bm{\theta}_1}\right\}^{-1}\bm{M}_{adj,1}(\bm{\theta}_1^*)+ o(1).
\end{align*}

We need the following conditions to make sure that $\lim\limits_{n\to \infty}\widehat{\gamma}_1 - \gamma_1^* = 0$ and subsequently $\lim\limits_{n\to \infty}\bm{M}_{adj,1}(\bm{\theta}_1^*) = \bm{0}$.

Condition 1. The kernel $K(\bm{u}_1)$ has bounded derivatives of order $\delta$, satisfies $\int K(\bm{u}_1) \mathrm{d} \bm{u}_1=1$, and has zero moments of order up to $m-1$ and nonzero $m$ th-order moment.

Condition 2. The true function of $s_1(\bm{u}_1)$ is continuously differentiable and bounded on an open set containing the support of $\bm{u}_1$.

Condition 3. The moment $E\{\exp (4 \gamma_1 \widehat{y}_{pse,1})\}$ is finite and the function $E\{\exp (4 \gamma_1 \widehat{y}_{pse,1}) \mid \bm{u}_1\} \pr(\bm{u}_1)$ is bounded, where $\pr(\bm{u}_1)$ is the marginal density of $\bm{u}_1$.

Condition 4. The bandwidth $c_1=c_{n_1}$ is such that $c_{n_1} \rightarrow 0, n_1 c_{n_1}^{p_1} \rightarrow \infty, n_1^{1 / 2} c_{n_1}^{p_1+2 \delta} / \log n_1 \rightarrow \infty$, and $n_1 c_{n_1}^{2 m} \rightarrow 0$ as $n_1 \rightarrow \infty$, where $p_1$ is the dimension of $\bm{u}_1$, $n_1$ is the number of units with fully-observed covariates at stage 1.

Let $\bm{B}_1^{\prime}(\bm{\theta}_2,\gamma_1) = \frac{\partial \bm{B}_1(\bm{\theta}_2,\gamma_1)^{\mathrm{T}}W_1\bm{B}_1(\bm{\theta}_2,\gamma_1)}{\partial \gamma_1}$, with conditions 1-4 and the convergence of $\bm{\theta}_2$, we have $\lim\limits_{n\to \infty}E\{\bm{B}_1^{\prime}(\widehat{\bm{\theta}}_2,\gamma_1^*)\}\stackrel{p}{\longrightarrow}0$ and $\lim\limits_{n\to \infty} \bm{W}_1 \stackrel{p}{\longrightarrow} \bm{W}_1^*$. 

By Taylor expansion, we have:
\begin{align*}
    &\left(\begin{array}{l}\widehat{\bm{\theta}}_2\\\widehat{\gamma}_1\end{array}\right) - \left(\begin{array}{l}
    \bm{\theta}_2^* \\ \gamma_1^* \end{array}\right) = - E\left[ \frac{\partial\{\bm{M}_2(\bm{\theta}_2^*),\bm{B}_1^{\prime}(\bm{\theta}_2^*,\gamma_1^*)\}}{\partial (\bm{\theta}_2, \gamma_1)}  \right]^{-1} \left(\begin{array}{l}
    \bm{M}_2(\bm{\theta}_2^*) \\ \bm{B}_1^{\prime}(\bm{\theta}_2^*,\gamma_1^*) \end{array}\right) + o(1)\\
    &= -\left( E\left[ \begin{array}{ll}
    \frac{\partial \bm{M}_2(\bm{\theta}_2^*)}{\partial \bm{\theta}_2} & \frac{\partial \bm{M}_2(\bm{\theta}_2^*)}{\partial \gamma_1}\\ \frac{\partial \bm{B}_1^{\prime}(\bm{\theta}_2^*,\gamma_1^*)}{\partial \bm{\theta}_2} & \frac{\partial \bm{B}_1^{\prime}(\bm{\theta}_2^*,\gamma_1^*)}{\partial \gamma_1}\end{array} \right]\right)^{-1}\left(\begin{array}{l}
    \bm{M}_2(\bm{\theta}_2^*) \\ \bm{B}_1^{\prime}(\bm{\theta}_2^*,\gamma_1^*) \end{array}\right) + o(1)\\
    & = -\left(\left[ \begin{array}{ll}
    E\left\{\frac{\partial \bm{M}_2(\bm{\theta}_2^*)}{\partial \bm{\theta}_2}\right\} & \bm{0}\\ E\left\{\frac{\partial \bm{B}_1^{\prime}(\bm{\theta}_2^*,\gamma_1^*)}{\partial \bm{\theta}_2}\right\} & E\left\{\frac{\partial \bm{B}_1^{\prime}(\bm{\theta}_2^*,\gamma_1^*)}{\partial \gamma_1}\right\} \end{array} \right]\right)^{-1} \left(\begin{array}{l}
    \bm{M}_2(\bm{\theta}_2^*) \\ \bm{B}_1^{\prime}(\bm{\theta}_2^*,\gamma_1^*) \end{array}\right) + o(1)\\
    &= \left[ \begin{array}{ll}
    -\left[E\left\{\frac{\partial \bm{M}_2(\bm{\theta}_2^*)}{\partial \bm{\theta}_2}\right\}\right]^{-1} & \bm{0}\\ \left[E\left\{ \frac{\partial \bm{B}_1^{\prime}(\bm{\theta}_2^*,\gamma_1^*)}{\partial \gamma_1}\right\}\right]^{-1} E\left\{\frac{\partial \bm{B}_1^{\prime}(\bm{\theta}_2^*,\gamma_1^*)}{\partial \bm{\theta}_2}\right\} \left[E\left\{\frac{\partial \bm{M}_2(\bm{\theta}_2^*)}{\partial \bm{\theta}_2}\right\}\right]^{-1} & -\left[E\left\{\frac{\partial \bm{B}_1^{\prime}(\bm{\theta}_2^*,\gamma_1^*)}{\partial \gamma_1}\right\}\right]^{-1} \end{array} \right] \left(\begin{array}{l} \bm{M}_2(\bm{\theta}_2^*) \\ \bm{B}_1^{\prime}(\bm{\theta}_2^*,\gamma_1^*) \end{array}\right) + o(1).
\end{align*}
Therefore, $\widehat{\gamma}_1 - \gamma_1^* = E\left\{\frac{\partial \bm{B}_1^{\prime}(\bm{\theta}_2^*,\gamma_1^*)}{\partial \bm{\theta}_2}\right\} \left[E\left\{\frac{\partial \bm{M}_2(\bm{\theta}_2^*)}{\partial \bm{\theta}_2}\right\}\right]^{-1}\bm{M}_2(\bm{\theta}_2^*) - \left[E\left\{\frac{\partial \bm{B}_1^{\prime}(\bm{\theta}_2^*,\gamma_1^*)}{\partial \gamma_1}\right\}\right]^{-1} \bm{B}_1^{\prime}(\bm{\theta}_2^*,\gamma_1^*)$,
\begin{align*}
    &M_{\mathrm{adj}, 1}(\bm{\theta}_1) = \bm{M}_1(\bm{\theta}_1, \bm{\theta}_2^*,\gamma_1^*) - E\left\{\frac{\partial \bm{M}_1(\bm{\theta}_1, \bm{\theta}_2^*,\gamma_1^*)}{\partial \bm{\theta}_2}\right\} \left[E\left\{\frac{\partial \bm{M}_2(\bm{\theta}_2^*)}{\partial \bm{\theta}_2}\right\}\right]^{-1} \bm{M}_2(\bm{\theta}_2^*)\\
    &+E\left\{\frac{\partial \bm{M}_1(\bm{\theta}_1, \bm{\theta}_2^*,\gamma_1^*)}{\partial \gamma_1}\right\}\left[E\left\{\frac{\partial \bm{B}_1^{\prime}(\bm{\theta}_2^*,\gamma_1^*)}{\partial \gamma_1}\right\}\right]^{-1}E\left\{\frac{\partial \bm{B}_1^{\prime}(\bm{\theta}_2^*,\gamma_1^*)}{\partial \bm{\theta}_2}\right\} \left[E\left\{\frac{\partial \bm{M}_2(\bm{\theta}_2^*)}{\partial \bm{\theta}_2}\right\}\right]^{-1}\bm{M}_2(\bm{\theta}_2^*)\\
    &-E\left\{\frac{\partial \bm{M}_1(\bm{\theta}_1, \bm{\theta}_2^*,\gamma_1^*)}{\partial \gamma_1}\right\}\left[E\left\{\frac{\partial \bm{B}_1^{\prime}(\bm{\theta}_2^*,\gamma_1^*)}{\partial \gamma_1} \right\} \right]^{-1} \bm{B}_1^{\prime}(\bm{\theta}_2^*,\gamma_1^*)
     + o(1).
\end{align*}


And we have $\widehat{\bm{\theta}}_{1,EE} - \bm{\theta}_1^* \rightarrow N(\bm{0},\Sigma_{\bm{\theta}_1} )$, where
$\Sigma_{\bm{\theta}_1} = E\left\{\left(\left[E\left\{\frac{\partial M_{adj,1}(\bm{\theta}_1^*)}{\partial \bm{\theta}_1} \right\}\right]^{-1} M_{adj,1}(\bm{\theta}_1^*) \right)^{\otimes 2}\right\} $.

%\iffalse

The detailed form of $\frac{\partial \bm{M}_2(\bm{\theta}_2^*)}{\partial \bm{\theta}_2}$ is provided in equation \eqref{eq1} in Web Appendix A. Assume that  $\frac{\partial \bm{B}_1^{\prime}(\bm{\theta}_2^*,\gamma_1^*)}{\partial \bm{\theta}_2}$ and $\frac{\partial \bm{B}_1^{\prime}(\bm{\theta}_2^*,\gamma_1^*)}{\partial \gamma_1}$ exist. To obtain $M_{\mathrm{adj}, 1}(\bm{\theta}_1)$ from observed data, we give the detailed form of $E\left\{\frac{\partial \bm{M}_1(\bm{\theta}_1, \bm{\theta}_2^*,\gamma_1^*)}{\partial \bm{\theta}_2}\right\}$, $E\left\{\frac{\partial \bm{M}_1(\bm{\theta}_1, \bm{\theta}_2^*,\gamma_1^*)}{\partial \gamma_1}\right\}$, $E\left\{\frac{\partial \bm{B}_1^{\prime}(\bm{\theta}_2^*,\gamma_1^*)}{\partial \bm{\theta}_2}\right\}$, and $E\left\{\frac{\partial \bm{B}_1^{\prime}(\bm{\theta}_2^*,\gamma_1^*)}{\partial \gamma_1}\right\}$ in the following text. Since the missing indicator of $Y_{pse,1}$ is mainly determined by the parametric form $Q_{2}$, under the assumption that the variables employed in the computation of $Y_{pse,1}$ remain unchanged, we have $\widehat{r}_{pse,1} = r_{pse,1}$. 

Since $\widehat{y}_{pse,1}$ is calculated with $\bm{\theta}_2$, we denote it as $\widehat{y}_{pse,1}(\widehat{\bm{\theta}}_2)$. Note that $\widehat{y}_{pse,1}(\bm{\theta}_2^*) = y_{pse,1}$. Consequently,
\begin{align*}
    &\frac{\partial \bm{M}_1(\bm{\theta}_1, \bm{\theta}_2^*,\gamma_1^*)}{\partial \bm{\theta}_2}  = \frac{\partial \bm{M}_1(\bm{\theta}_1, \widehat{y}_{pse,1}(\bm{\theta}_2^*),\gamma_1^*)}{\partial \left\{\widehat{y}_{pse,1}(\bm{\theta}_2^*)\right\}} \frac{\partial \left\{\widehat{y}_{pse,1}(\bm{\theta}_2^*)\right\}}{\partial \bm{\theta}_2}\\
    &= \frac{\partial \bm{M}_1(\bm{\theta}_1, y_{pse,1},\gamma_1^*)}{\partial y_{pse,1}} \frac{\partial y_{pse,1}}{\partial \bm{\theta}_2}\\
    &= \frac{\partial \widehat{E} \left(\frac{r_{pse,1}}{\widehat{\pi}_{1}}r_{1} \frac{\partial \left[ \left\{ Q_1(\bm{h}_1,a_1;\bm{\theta}_1) - y_{pse,1}  \right\}^2\right]}{\partial \bm{\theta}_1} \right)}{\partial y_{pse,1}} \frac{\partial y_{pse,1}}{\partial \bm{\theta}_2}\\
    &= \widehat{E} \left( 2 r_{pse,1} r_1 \frac{\partial Q_1(\bm{h}_1,a_1;\bm{\theta}_1)}{\partial \bm{\theta}_1} \frac{\partial [ \widehat{\pi}_1\{ Q_1(\bm{h}_1,a_1;\bm{\theta}_1) - y_{pse,1}  \}] }{\partial y_{pse,1}} \right) \frac{\partial y_{pse,1}}{\partial \bm{\theta}_2}\\
    &= \widehat{E} \left( 2 r_{pse,1} r_1 \frac{\partial Q_1(\bm{h}_1,a_1;\bm{\theta}_1)}{\partial \bm{\theta}_1} \frac{\partial [ 1/\widehat{\pi}_1\{ Q_1(\bm{h}_1,a_1;\bm{\theta}_1) - y_{pse,1}  \}] }{\partial y_{pse,1}} \right) \frac{\partial y_{pse,1}}{\partial \bm{\theta}_2}\\  
    &= \widehat{E} \left( 2 r_{pse,1} r_1 \frac{\partial Q_1(\bm{h}_1,a_1;\bm{\theta}_1)}{\partial \bm{\theta}_1} \left[ - \frac{1}{\widehat{\pi}_1} + \{ Q_1(\bm{h}_1,a_1;\bm{\theta}_1) - y_{pse,1}\} \frac{\partial ( 1/\widehat{\pi}_1)}{\partial y_{pse,1}} \right]\right) \frac{\partial y_{pse,1}}{\partial \bm{\theta}_2}\\
    &= \widehat{E} \left( 2 r_{pse,1} r_1 \frac{\partial Q_1(\bm{h}_1,a_1;\bm{\theta}_1)}{\partial \bm{\theta}_1} \left[ \{ \gamma_1^* Q_1(\bm{h}_1,a_1;\bm{\theta}_1) - \gamma_1^* y_{pse,1}-1\} \exp\{\widehat{s}_{\gamma_1^*}(\bm{u}_1) + \gamma_1^*y_{pse,1}\} -1 \right]\right) \frac{\partial y_{pse,1}}{\partial \bm{\theta}_2},
\end{align*}
where $\frac{\partial Q_1(\bm{h}_1,a_1;\bm{\theta}_1)}{\partial \bm{\theta}_1}$ and $\frac{\partial y_{pse,1}}{\partial \bm{\theta}_2}$ depends on the specified form of the Q-functions.

\iffalse
We derive the detailed form of $\frac{\partial ( 1/\widehat{\pi}_1)}{\partial y_{pse,1}}$ as follows
\begin{align}
    &\frac{\partial  ( 1/\widehat{\pi}_1)}{\partial y_{pse,1}} =
    \frac{\partial \exp\{\widehat{s}_{\gamma_1^*}(\bm{u}_1) \}}{\partial y_{pse,1}} \exp(\gamma_1^*y_{pse,1}) +  \frac{\partial \exp(\gamma_1^*y_{pse,1})}{\partial y_{pse,1}}\exp\{\widehat{s}_{\gamma_1^*}(\bm{u}_1) \} \notag \\
    &= \frac{\partial \left\{\frac{\sum_{i=1}^{n} \left(1-r_{pse,1,i}\right)  K_{c_1}\left(\bm{u}_1-\bm{u}_{1,i}\right) r_{1,i} }{\sum_{i=1}^{n}  r_{pse,1,i}   \exp \left(\gamma_1^*y_{pse,1,i}\right) K_{c_1}\left(\bm{u}_1-\bm{u}_{1,i}\right) r_{1,i}} \right\}}{\partial y_{pse,1}} \exp(\gamma_1^*y_{pse,1}) +   \gamma_1^* \exp\{\widehat{s}_{\gamma_1^*}(\bm{u}_1) + \gamma_1^*y_{pse,1}\}  \notag \\
    &= \frac{ [\sum_{i=1}^{n} \{\left(1-r_{pse,1,i}\right)  K_{c_1}\left(\bm{u}_1-\bm{u}_{1,i}\right) r_{1,i}\} ] \sum_{i=1}^{n} -\gamma_1^* r_{pse,1,i}   \exp \left(\gamma_1^*y_{pse,1,i}\right) K_{c_1}\left(\bm{u}_1-\bm{u}_{1,i}\right) r_{1,i}}{(\sum_{i=1}^{n}  r_{pse,1,i}   \exp \left(\gamma_1^*y_{pse,1,i}\right) K_{c_1}\left(\bm{u}_1-\bm{u}_{1,i}\right) r_{1,i})^2}  \exp(\gamma_1^*y_{pse,1})  \notag  \\
    &\qquad  +   \gamma_1^* \exp\{\widehat{s}_{\gamma_1^*}(\bm{u}_1) + \gamma_1^*y_{pse,1}\}  \notag \\
    & = -\gamma_1^* \frac{\sum_{i=1}^{n} \left(1-r_{pse,1,i}\right)  K_{c_1}\left(\bm{u}_1-\bm{u}_{1,i}\right) r_{1,i} }{\sum_{i=1}^{n}  r_{pse,1,i}   \exp \left(\gamma_1^*y_{pse,1,i}\right) K_{c_1}\left(\bm{u}_1-\bm{u}_{1,i}\right) r_{1,i}} \exp(\gamma_1^*y_{pse,1}) +   \gamma_1^* \exp\{\widehat{s}_{\gamma_1^*}(\bm{u}_1) + \gamma_1^*y_{pse,1}\}  \notag \\
    &= -\gamma_1^* \exp\{\widehat{s}_{\gamma_1^*}(\bm{u}_1)\} \exp(\gamma_1^*y_{pse,1}) +   \gamma_1^* \exp\{\widehat{s}_{\gamma_1^*}(\bm{u}_1) + \gamma_1^*y_{pse,1}\}  \notag  \\
    &=0.
\end{align}
Therefore, $\frac{\partial \bm{M}_1(\bm{\theta}_1, \bm{\theta}_2^*,\gamma_1^*)}{\partial \bm{\theta}_2} = \widehat{E} \left\{ \frac{2 r_{pse,1} r_1}{\widehat{\pi}_1} \frac{\partial Q_1(\bm{h}_1,a_1;\bm{\theta}_1)}{\partial \bm{\theta}_1}   \right\} \frac{\partial y_{pse,1}}{\partial \bm{\theta}_2}$. The detailed form of $\frac{\partial Q_1(\bm{h}_1,a_1;\bm{\theta}_1)}{\partial \bm{\theta}_1}$ and $\frac{\partial y_{pse,1}}{\partial \bm{\theta}_2}$ depends on the specified form of the Q-functions.

Similarly, the influence of $\widehat{\bm{\theta}}_2$ on $\bm{B}_1'$ is solely mediated through its impact on $\widehat{Y}_{pse,1}$. And the influence of $\widehat{Y}_{pse,1}$ on $\bm{B}_1'$ is solely mediated through its impact on $\widehat{\pi}_{1}$. Thus, 
\begin{align*}
    \frac{\partial \bm{B}_1'(\bm{\theta}_2^*, \gamma_1^*)}{\partial \bm{\theta}_2} &= \frac{\partial \bm{B}_1'(\bm{\theta}_2^*, \gamma_1^*)}{\partial \widehat{\pi}_1(\bm{u}_1,\widehat{y}_{pse,1}(\bm{\theta}_2^*),\gamma_1^*)} \frac{\partial\widehat{\pi}_1(\bm{u}_1,\widehat{y}_{pse,1}(\bm{\theta}_2^*),\gamma_1^*)}{\partial \widehat{y}_{pse,1}(\bm{\theta}_2^*)}  \frac{\partial \widehat{y}_{pse,1}(\bm{\theta}_2^*)}{\partial\bm{\theta}_2}\\
    &=\frac{\partial \bm{B}_1'(\bm{\theta}_2^*, \gamma_1^*)}{\partial \widehat{\pi}_1(\bm{u}_1,y_{pse,1},\gamma_1^*)} \frac{\partial\widehat{\pi}_1(\bm{u}_1,y_{pse,1},\gamma_1^*)}{\partial y_{pse,1}}  \frac{\partial y_{pse,1}}{\partial\bm{\theta}_2} \label{eq1}
\end{align*} 

As we derived in equation \eqref{eq1}, $\frac{\partial\widehat{\pi}_1(\bm{u}_1,y_{pse,1},\gamma_1^*)}{\partial y_{pse,1}} = 0$. Therefore, $\frac{\partial \bm{B}_1'(\bm{\theta}_2^*, \gamma_1^*)}{\partial \bm{\theta}_2} = 0$.
\fi
Similarly, we obtain $\frac{\partial \bm{M}_1(\bm{\theta}_1, \bm{\theta}_2^*,\gamma_1^*)}{\partial \gamma_1}$ using the following equations
\begin{align}
    &\frac{\partial \bm{M}_1(\bm{\theta}_1, \bm{\theta}_2^*,\gamma_1^*)}{\partial \gamma_1}  = \frac{\partial \bm{M}_1\{\bm{\theta}_1, \widehat{y}_{pse,1}(\bm{\theta}_2^*),\gamma_1^*\}}{\partial \left[1/\widehat{\pi}_1\{\bm{u}_1,\widehat{y}_{pse,1}(\bm{\theta}_2^*),\gamma_1^*\}\right]} \frac{\partial \left[1/\widehat{\pi}_1\{\bm{u}_1,\widehat{y}_{pse,1}(\bm{\theta}_2^*),\gamma_1^*\}\right]}{\partial \gamma_1} \notag \\
    &=  \frac{\partial \bm{M}_1(\bm{\theta}_1, y_{pse,1},\gamma_1^*)}{\partial \left\{1/\widehat{\pi}_1(\bm{u}_1,y_{pse,1},\gamma_1^*)\right\}} \frac{\partial \left\{1/\widehat{\pi}_1(\bm{u}_1,y_{pse,1},\gamma_1^*)\right\}}{\partial \gamma_1} \notag \\
    &= \frac{\partial \widehat{E} \left(\frac{r_{pse,1}}{\widehat{\pi}_{1}}r_{1} \frac{\partial \left[ \left\{ Q_1(\bm{h}_1,a_1;\bm{\theta}_1) - y_{pse,1}  \right\}^2\right]}{\partial \bm{\theta}_1} \right)}{\partial (1/\widehat{\pi}_1)} \frac{\partial \left\{1/\widehat{\pi}_1(\bm{u}_1,y_{pse,1},\gamma_1^*)\right\}}{\partial \gamma_1} \notag \\
    &= \widehat{E} \bigg[2r_{pse,1}r_{1}\left\{ Q_1(\bm{h}_1,a_1;\bm{\theta}_1) - y_{pse,1}\right\} \frac{\partial  Q_1(\bm{h}_1,a_1;\bm{\theta}_1) }{\partial \bm{\theta}_1} \bigg] \bigg[ \frac{\partial \left\{1/\widehat{\pi}_1(\bm{u}_1,y_{pse,1},\gamma_1^*)\right\}}{\partial \gamma_1}\bigg]. \label{eq2}
\end{align}

The explicit expression for $\frac{\partial ( 1/\widehat{\pi}_1)}{\partial \gamma_{1}}$ can be derived with the following equations
\begin{align}
    &\frac{\partial \{ 1/\widehat{\pi}_1(\bm{u}_1,y_{pse,1},\gamma_1^*)\}}{\partial \gamma_{1}} = \frac{\partial \left[1+\exp\{\widehat{s}_{\gamma_1^*}(\bm{u}_1)\}\exp(\gamma_1^*y_{pse,1})\right]}{\partial \gamma_1} \notag \\
    & = \gamma_1^*\exp\{\widehat{s}_{\gamma_1^*}(\bm{u}_1)\}\exp(\gamma_1^*y_{pse,1}) + \frac{\partial \exp\{\widehat{s}_{\gamma_1^*}(\bm{u}_1)\}}{\partial \gamma_1} \exp(\gamma_1^*y_{pse,1}) \notag \\
    & = \gamma_1^*\exp\{\widehat{s}_{\gamma_1^*}(\bm{u}_1)\}\exp(\gamma_1^*y_{pse,1}) + \frac{\partial \left\{\frac{\sum_{i=1}^{n} \left(1-r_{pse,1,i}\right)  K_{c_1}\left(\bm{u}_1-\bm{u}_{1,i}\right) r_{1,i} }{\sum_{i=1}^{n}  r_{pse,1,i}   \exp \left(\gamma_1^* y_{pse,1,i}\right) K_{c_1}\left(\bm{u}_1-\bm{u}_{1,i}\right) r_{1,i}} \right\}}{\partial \gamma_1} \exp(\gamma_1^*y_{pse,1}) \notag \\
    & =\exp\{\widehat{s}_{\gamma_1^*}(\bm{u}_1)\}\exp(\gamma_1^*y_{pse,1}) \left\{\gamma_1^*+ \frac{\sum_{i=1}^{n} r_{pse,1,i} y_{pse,1,i}  \exp \left(\gamma_1^* y_{pse,1,i}\right) K_{c_1}\left(\bm{u}_1-\bm{u}_{1,i}\right) r_{1,i} }{\sum_{i=1}^{n}  r_{pse,1,i}   \exp \left(\gamma_1^* y_{pse,1,i}\right) K_{c_1}\left(\bm{u}_1-\bm{u}_{1,i}\right) r_{1,i}}\right\}. \label{eq3}
\end{align}

Combining equations \eqref{eq2} and \eqref{eq3}, we can obtain the detailed form of $\frac{\partial \bm{M}_1(\bm{\theta}_1, \bm{\theta}_2^*,\gamma_1^*)}{\partial \gamma_1}$.


%\fi

The detailed form of $E\{\frac{\partial \bm{B}_1^{\prime}(\bm{\theta}_2^*,\gamma_1^*)}{\partial \gamma_1}\}$ and $E\{\frac{\partial \bm{B}_1^{\prime}(\bm{\theta}_2^*,\gamma_1^*)}{\partial \bm{\theta}_2}\}$ can be derived following \cite{Shao_Wang_2016} when the following regularity conditions hold.


Condition 5. There is a vector of the functional $G(r_{pse,1}, \bm{u}_1, \bm{\omega})$ which is linear in $\bm{\omega}=\left(\omega_{1}, \omega_{2}\right)^{\mathrm{T}}$ and such that:

(i) for small enough $\left\|\bm{\omega}-\bm{\omega}^*\right\|$, $\| \widetilde{m}\left(r_{pse,1}, \bm{u}_1, \widehat{y}_{pse,1}, \bm{\omega}, \gamma_1\right)-\widetilde{m}\left(r_{pse,1}, \bm{u}_1, \widehat{y}_{pse,1}, \bm{\omega}^*, \gamma_1\right)-G(r_{pse,1},\\ \bm{u}_1, \bm{\omega}-\bm{\omega}^*) \| \leqslant c(r_{pse,1}, \bm{u}_1)\left(\left\|\bm{\omega}-\bm{\omega}^*\right\|\right)^2$, where $\widetilde{m}(r_{pse,1}, \bm{u}_1, \widehat{y}_{pse,1}, \bm{\omega}, \gamma_1)$ is the $L$-dimensional vector. And its $l$-th component $\widetilde{m}_l(r_{pse,1}, \bm{u}_1, \widehat{y}_{pse,1}, \bm{\omega}, \gamma_1)$, is defined as $\mathbb{I}(\bm{z}_1=l)[r_{pse,1}\{1+\exp (\gamma_1 \widehat{y}_{pse,1}) \omega_1(\bm{u}_1) / \omega_2(\bm{u}_1)\}-1]$, where $\bm{\omega}^* = [E(1-r_{pse,1} \mid \bm{u}_1), E(r_{pse,1} \exp (\gamma_1^* y_{pse,1}) \mid \bm{u}_1)]^{\mathrm{T}}$, and $E\{c(r_{pse,1}, \bm{u}_1)\}<\infty$;
%$\omega_1 = E[1-r_{pse,1} \mid \bm{u}_1]$, $\omega_2 = E[r_{pse,1}\\ \exp \left(\gamma_1 y_{pse,1}\right) \mid \bm{u}_1]$, and $E[c(r_{pse,1}, \bm{u}_1, y_{pse,1})]<\infty$;

(ii) $\|G(r_{pse,1}, \bm{u}_1, \bm{\omega})\| \leqslant e(r_{pse,1}, \bm{u}_1)\| \bm{\omega}\|$ and $E\{e(r_{pse,1}, \bm{u}_1)^2\}<\infty$;

(iii) there exists an almost everywhere continuous function $v_1(\bm{u}_1)$ with $\int\|v_1(\bm{u}_1)\| \mathrm{d} \bm{u}_1<\infty$, $E\{G(r_{pse,1}, \bm{u}_1, \bm{\omega})\}=\int v_1(\bm{u}_1) \bm{\omega}(\bm{u}_1) \mathrm{d} \bm{u}_1$ for all $\|\bm{\omega}\| \leqslant \infty$, and $E\{\sup _{\|\zeta\| \leqslant \epsilon}\|v_1(\bm{u}_1+\zeta)\|^4\}<\infty$ for some $\epsilon>0$.

Condition 6. For small enough $\left\|\bm{\omega}-\bm{\omega}_0\right\|, \widetilde{m}(r_{pse,1}, \bm{u}_1, \widehat{y}_{pse,1}, \bm{\omega}, \gamma_1)$ is continuously differentiable in $\gamma_1$ in a neighbourhood of $(y_{pse,1},\gamma_{1}^*)$. And there is a function $k(r_{pse,1}, \bm{u}_1)$ satisfying $\left\| \frac{\partial \widetilde{m}(r_{pse,1}, \bm{u}_1, \widehat{y}_{pse,1}, \bm{\omega}, \gamma_1)}{\partial \gamma_1}- \frac{\partial \widetilde{m}\left(r_{pse,1}, \bm{u}_1, y_{pse,1}, \bm{\omega}^*, \gamma_{1}^*\right)}{\partial \gamma_1}\right\| \leq k(r_{pse,1}, \bm{u}_1)(\left\|\gamma_{1}^*-\gamma_1\right\|^\epsilon+\left\|\bm{\omega}-\bm{\omega}^*\right\|^\epsilon+\left\|\widehat{y}_{pse,1}-y_{pse,1}\right\|^\epsilon)$ for an $\epsilon>0$, $E\{k(r_{pse,1}, \bm{u}_1)\}<\infty$. Besides, $  \bm{\Gamma}_{\gamma_1} = E\left\{\frac{\partial \widetilde{m}\left(r_{pse,1}, \bm{u}_1, y_{pse,1}, \bm{\omega}^*, \gamma_{1}^*\right)}{\partial \gamma_1}\right\}$ and $\bm{\Gamma}_{y_{pse,1}} = E\left\{\frac{\partial \widetilde{m}\left(r_{pse,1}, \bm{u}_1, y_{pse,1}, \bm{\omega}^*, \gamma_{1}^*\right)}{\partial \widehat{y}_{pse,1}}\right\}$ exist and are of full rank.


Assume that Conditions $1-6$ hold. Then, as $n \rightarrow \infty$, $E\left\{\frac{\partial \bm{B}_1'(\bm{\theta}_2^*,\gamma_1^*)}{\partial \gamma_1}\right\} \to 2\bm{\Gamma}_{\gamma_1}^{\mathrm{T}}\bm{W}_1^*\bm{\Gamma}_{\gamma_1}$, $E\left\{\frac{\partial \bm{B}_1'(\bm{\theta}_2^*,\gamma_1^*)}{\partial \bm{\theta}_2}\right\} \to 2\bm{\Gamma}_{y_{pse,1}}^{\mathrm{T}}\bm{W}_1^*\bm{\Gamma}_{y_{pse,1}}\frac{\partial y_{pse,1}}{\partial \bm{\theta}_2}$. 
\iffalse
Besides, $\bm{B}_1^{\prime}(\bm{\theta}_2^*,\gamma_1^*) \to  2\bm{\Gamma}_1^{\mathrm{T}} \bm{W}_1 \left[  m\left(r_{pse,1}, \bm{u}_1, y_{pse,1}, s_{\gamma_1}^*, \gamma_1^*\right)+ \tau_1\left(r_{pse,1}, \bm{u}_1, y_{pse,1}, \gamma_1^*\right) \right]$, where
$\tau_1\left(r_{pse,1}, \bm{u}_1, y_{pse,1}, \gamma_1^*\right)=v_1(\bm{u}_1)\left\{\begin{array}{c}
1-r_{pse,1} \\
r_{pse,1} \exp \left(\gamma_1^* y_{pse,1}\right)
\end{array}\right\}-E\left[v_1(\bm{u}_1)\left\{\begin{array}{c}
1-r_{pse,1} \\
r_{pse,1} \exp \left(\gamma_1^* y_{pse,1}\right)
\end{array}\right\}\right].$
\fi

As mentioned by \cite{Shao_Wang_2016}, the form of the asymptotic variance of this semi-parameteric estimator is complicated. Consequently, we suggest using the bootstrap method to estimate it.

Besides, according to the Law of Large Numbers, $\widehat{\bm{\theta}}_{1,EE}$ converges to $\bm{\theta}_1$ in probability. Therefore, following the proof of Theorem 1, we can similarly prove the convergence in probability of $\widehat{d}^{\mathrm{opt}}_1$ by proving that $Q_1(\bm{h}_1,a_1;\widehat{\bm{\theta}}_{1,EE}) \stackrel{p}{\longrightarrow} Q_1(\bm{h}_1,a_1;\bm{\theta}_{1})$ for any fixed $(\bm{h}_1,a_1)$.  


\section{Web Appendix E: Proof of Theorem 3}

We need the following conditions to guarantee that $\lim\limits_{n\to \infty}\exp\{\widehat{s}'_{\gamma'_1,1}(\bm{g}_1) \} = \exp\{s'_{\gamma'_1,1}(\bm{g}_1) \}$:

Condition 7. The kernel $K'(\bm{g}_1)$ has bounded derivatives of order $\delta'$, satisfies \\ $\int K'(\bm{g}_1) \mathrm{d} \bm{g}_1=1$, and has zero moments of order up to $m-1$ and nonzero $m$ th-order moment.

Condition 8. The true function of $s_1'(\bm{g}_1)$ is continuously differentiable and bounded on an open set containing the support of $\bm{g}_1$.

Condition 9. The moment $E[\exp (4 \gamma_1 \widehat{y}_{pse,1})]$ is finite and the function $E[\exp (4 \gamma_1 \widehat{y}_{pse,1}) \mid \bm{g}_1] \pr(\bm{g}_1)$ is bounded, where $\pr(\bm{g}_1)$ is the marginal density of $\bm{g}_1$.

Condition 10. The bandwidth $c'_1=c'_{n_1}$ is such that $c'_{n_1} \rightarrow 0, n_1 {c'_{n_1}}^{p_1'} \rightarrow \infty, n_1^{1/2} {c'_{n_1}}^{p_1'+2 \delta'} / \log n_1 \rightarrow \infty$, and $n_1 {c'_{n_1}}^{2 m} \rightarrow 0$ as $n_1 \rightarrow \infty$, where $p_1'$ is the dimension of $\bm{g}_1$, $n_1$ is the number of units with fully-observed covariates at stage 1.

Under these conditions, $\exp\{\widehat{s}'_{\gamma'_1,1}(\bm{h}_1)\}$ is a consistent estimator of $\exp\{s'_{\gamma'_1,1}(\bm{h}_1)\}$. Therefore, the estimated missingness probability model is consistent. And $\bm{M}_1(\bm{\theta}_1)$ is an unbiased estimation equation for $\bm{\theta}_1$ \citep{Newey_McFadden_1994}. Similar to Web Appendix D, where we provided the asymptotic properties of $\widehat{\bm{\theta}}_1$ for the  WQ-EE method, we establish the asymptotic normality of $\widehat{\bm{\theta}}_1$ for the WQ-SA method under the regularity condition that $P(\{q_{2,1}(\bm{H}_2; \widehat{\bm{\psi}}_2) = 0\})=0$. 

Let  $
M_{\mathrm{adj}, 1}'(\bm{\theta}_1) = \bm{M}_1(\bm{\theta}_1, \bm{\theta}_2^*) + E\left\{ \frac{\partial \bm{M}_1(\bm{\theta}_1, \bm{\theta}_2^*)}{\partial \bm{\theta}_2} \right\}
(\widehat{\bm{\theta}}_2 -\bm{\theta}_2^*) + o(1)
$. Then we have
\begin{align*}
    M_{\mathrm{adj}, 1}'(\bm{\theta}_1) = \bm{M}_1(\bm{\theta}_1, \bm{\theta}_2^*) - E\left\{ \frac{\partial \bm{M}_1(\bm{\theta}_1, \bm{\theta}_2^*)}{\partial \bm{\theta}_2} \right\}\left\{ \frac{\partial \bm{M}_2(\bm{\theta}_2^*)}{\partial \bm{\theta}_2} \right\}^{-1}\bm{M}_2(\bm{\theta}_2^*).
\end{align*}

\iffalse
\begin{align*}
    &E\left[ \frac{\partial \bm{M}_1(\bm{\theta}_1, \bm{\theta}_2)}{\partial \bm{\theta}_2} \right] = E\left[ \frac{\partial \left[\frac{R_{pse,1}}{\widehat{\pi}_{1}} \mathbb{I}(\bm{R}_{1}=1) \frac{\partial \left[ \left\{ Q_t(\bm{h}_t,a_t;\bm{\theta}_1) - Y_{pse,1}  \right\}^2\right]}{\partial (\bm{\theta}_1)} \right]}{\partial \bm{\theta}_2} \right]\\
    &\qquad = E\left[  \frac{ \partial \left[ 2\mathbb{I}(\overline{\bm{R}}_{2}=\bm{1}_2) \widehat{\pi}^{-1}\{ Q_1(\bm{h}_1,a_1;\bm{\theta}_1)-Y_{pse,1} \} \frac{\partial Q_1(\bm{h}_1,a_1;\bm{\theta}_1)}{\partial \bm{\theta}_1} \right] }{\partial \bm{\theta}_2} \right]\\
    &\qquad = E\left[ 2\mathbb{I}(\overline{\bm{R}}_{2}=\bm{1}_2)\frac{\partial Q_1(\bm{h}_1,a_1;\bm{\theta}_1)}{\partial \bm{\theta}_1} \frac{ \partial \left[ \widehat{\pi}^{-1}\{ Q_1(\bm{h}_1,a_1;\bm{\theta}_1)-Y_{pse,1} \}\right]  }{\partial \bm{\theta}_2} \right] \\
    &\qquad = E\left[ 2\mathbb{I}(\overline{\bm{R}}_{2}=\bm{1}_2)\frac{\partial Q_1(\bm{h}_1,a_1;\bm{\theta}_1)}{\partial \bm{\theta}_1} \{ Q_1(\bm{h}_1,a_1;\bm{\theta}_1)-Y_{pse,1} \}\frac{ \partial  \widehat{\pi}^{-1} }{\partial \bm{\theta}_2} \right]\\
    &\qquad \qquad - E\left[ 2\mathbb{I}(\overline{\bm{R}}_{2}=\bm{1}_2)\frac{\partial Q_1(\bm{h}_1,a_1;\bm{\theta}_1)}{\partial \bm{\theta}_1} \widehat{\pi}^{-1} \frac{\partial Y_{pse,1}}{\partial \bm{\theta}_2} \right]
\end{align*}

\begin{align*}
    \frac{ \partial  \widehat{\pi}^{-1} }{\partial \bm{\theta}_2} = \frac{\partial [1+\exp(\widehat{s}_{\gamma_1,1})\exp(\gamma_1Y_{pse,1})]}{\partial \bm{\theta}_2}
\end{align*}
\fi


And $\widehat{\bm{\theta}}_{1,SA} - \bm{\theta}_1 \rightarrow N(\bm{0},\Sigma'_{\bm{\theta}_1} )$, where
$\Sigma'_{\bm{\theta}_1} = E\left\{\left(\left[E\left\{\frac{\partial M'_{adj,1}(\bm{\theta}_1)}{\partial \bm{\theta}_1} \right\}\right]^{-1} M'_{adj,1}(\bm{\theta}_1) \right)^{\otimes 2}\right\} $. The explicit mathematical expression for $\frac{\partial \bm{M}_2(\bm{\theta}_2^*)}{\partial \bm{\theta}_2}$ can be found in equation \eqref{eq1} in Web Appendix C. Employing a similar approach to that outlined in Web Appendix D, we obtain $\frac{\partial M'_{adj,1}(\bm{\theta}_1)}{\partial \bm{\theta}_1}  =\\ \widehat{E} \left( 2 r_{pse,1} r_1 \frac{\partial Q_1(\bm{h}_1,a_1;\bm{\theta}_1)}{\partial \bm{\theta}_1} \left[ \{ \gamma_1 Q_1(\bm{h}_1,a_1;\bm{\theta}_1) - \gamma_1 y_{pse,1}-1\} \exp\{\widehat{s}_{\gamma_1^*}(\bm{h}_1) + \gamma_1 y_{pse,1}\} -1 \right]\right) \frac{\partial y_{pse,1}}{\partial \bm{\theta}_2}$, where $\frac{\partial Q_1(\bm{h}_1,a_1;\bm{\theta}_1)}{\partial \bm{\theta}_1}$ and $\frac{\partial y_{pse,1}}{\partial \bm{\theta}_2}$ depend on the specified form of the Q-functions. Because of the complexity of the asymptotic variance, we recommend utilizing the bootstrap method for its estimation. Following the derivation of Theorem 1, we can prove the convergence in probability of $\widehat{d}^{\mathrm{opt}}_1$ by proving that $Q_1(\bm{h}_1,a_1;\widehat{\bm{\theta}}_{1, SA}) \stackrel{p}{\longrightarrow} Q_1(\bm{h}_1,a_1;\bm{\theta}_{1})$ for any fixed $(\bm{h}_1,a_1)$.

\section{Web Appendix F: Proof of Lemma 1}

Using  Bayes' Theorem, we have
\begin{align*}
    \pr(y_{pse,t}\mid \bm{h}_t,a_t,\overline{\bm{R}}_t = \bm{1}_t,R_{pse,t}=0) &= \frac{\pr(y_{pse,t},\bm{h}_t,a_t,\overline{\bm{R}}_t = \bm{1}_t,R_{pse,t}=0)}{\pr (\bm{h}_t,a_t,\overline{\bm{R}}_t = \bm{1}_t,R_{pse,t}=0)},\\
    \pr(y_{pse,t}\mid \bm{h}_t,a_t,\overline{\bm{R}}_t = \bm{1}_t,R_{pse,t}=1) &= \frac{\pr(y_{pse,t},\bm{h}_t,a_t,\overline{\bm{R}}_t = \bm{1}_t,R_{pse,t}=1)}{\pr (\bm{h}_t,a_t,\overline{\bm{R}}_t = \bm{1}_t,R_{pse,t}=1)} .
\end{align*}

This leads to the following equations:
\begin{align*}
    &\frac{\pr(y_{pse,t}\mid \bm{h}_t,a_t,\overline{\bm{R}}_t = \bm{1}_t,R_{pse,t}=0)}{\pr(y_{pse,t}\mid \bm{h}_t,a_t,\overline{\bm{R}}_t = \bm{1}_t,R_{pse,t}=1)} \\
    &\quad = \frac{\pr(y_{pse,t},\bm{h}_t,a_t,\overline{\bm{R}}_t = \bm{1}_t,R_{pse,t}=0)}{\pr (\bm{h}_t,a_t,\overline{\bm{R}}_t = \bm{1}_t,R_{pse,t}=0)}\frac{\pr (\bm{h}_t,a_t,\overline{\bm{R}}_t = \bm{1}_t,R_{pse,t}=1)}{\pr(y_{pse,t},\bm{h}_t,a_t,\overline{\bm{R}}_t = \bm{1}_t,R_{pse,t}=1)}\\
    &\quad = \frac{P(R_{pse,t}=0\mid y_{pse,t},\bm{h}_t,a_t,\overline{\bm{R}}_t = \bm{1}_t)}{P(R_{pse,t}=1\mid y_{pse,t},\bm{h}_t,a_t,\overline{\bm{R}}_t = \bm{1}_t)}\frac{P(R_{pse,t}=1\mid \bm{h}_t,a_t,\overline{\bm{R}}_t = \bm{1}_t)}{P(R_{pse,t}=0 \mid \bm{h}_t,a_t,\overline{\bm{R}}_t = \bm{1}_t)}.
\end{align*}

Under Assumption 6, we have $\frac{P(R_{pse,t}=0\mid y_{pse,t},\bm{h}_t,a_t,\overline{\bm{R}}_t = \bm{1}_t)}{P(R_{pse,t}=1\mid y_{pse,t},\bm{h}_t,a_t,\overline{\bm{R}}_t = \bm{1}_t)} = \exp\{s'_t(\bm{h}_t,a_t)+\gamma'_t y_{pse,t}\}$. Besides, 
\begin{align*}
    &\frac{P(R_{pse,t}=0 \mid \bm{h}_t,a_t,\overline{\bm{R}}_t = \bm{1}_t)}{P(R_{pse,t}=1\mid \bm{h}_t,a_t,\overline{\bm{R}}_t = \bm{1}_t)} \\
    & \qquad = \int \pr(y_{pse,t} \mid \bm{h}_t,a_t,\overline{\bm{R}}_t = \bm{1}_t,R_{pse,t}=0) \frac{P(R_{pse,t}=0 \mid \bm{h}_t,a_t,\overline{\bm{R}}_t = \bm{1}_t)}{P(R_{pse,t}=1\mid \bm{h}_t,a_t,\overline{\bm{R}}_t = \bm{1}_t)} \mathrm{d} y_{pse,t}\\
    & \qquad =\int \frac{\pr\left(y_{pse,t} \mid \bm{h}_t,a_t,\overline{\bm{R}}_t = \bm{1}_t\right)}{P\left(R_{pse,t}=1 \mid \bm{h}_t,a_t,\overline{\bm{R}}_t = \bm{1}_t\right)} P\left(R_{pse,t}=0 \mid \bm{h}_t,a_t,\overline{\bm{R}}_t = \bm{1}_t, y_{pse,t}\right) \mathrm{d} y_{pse,t}\\
    & \qquad =\int \frac{\pr\left(y_{pse,t} \mid \bm{h}_t,a_t,\overline{\bm{R}}_t = \bm{1}_t, R_{pse,t}=1\right)}{P\left(R_{pse,t}=1 \mid \bm{h}_t,a_t,\overline{\bm{R}}_t = \bm{1}_t, y_{pse,t}\right)} P\left(R_{pse,t}=0 \mid \bm{h}_t,a_t,\overline{\bm{R}}_t = \bm{1}_t, y_{pse,t}\right) \mathrm{d} y_{pse,t}\\
    & \qquad =\int \frac{P\left(R_{pse,t}=0 \mid \bm{h}_t,a_t,\overline{\bm{R}}_t = \bm{1}_t, y_{pse,t}\right)}{P\left(R_{pse,t}=1 \mid \bm{h}_t,a_t,\overline{\bm{R}}_t = \bm{1}_t, y_{pse,t}\right)} \pr\left(y_{pse,t} \mid \bm{h}_t,a_t,\overline{\bm{R}}_t = \bm{1}_t, R_{pse,t}=1\right) \mathrm{d} y_{pse,t} \\
    & \qquad = \int \exp\{s'_t(\bm{h}_t,a_t)+\gamma'_t y_{pse,t}\} \pr\left(y_{pse,t} \mid \bm{h}_t,a_t,\overline{\bm{R}}_t = \bm{1}_t, R_{pse,t}=1\right) \mathrm{d} y_{pse,t}\\
    & \qquad = E[\exp\{s'_t(\bm{h}_t,a_t)+\gamma'_t y_{pse,t}\} \mid \bm{h}_t,a_t,\overline{\bm{R}}_t = \bm{1}_t, R_{pse,t}=1].
\end{align*}

Therefore, 
\begin{align*}
    \frac{\pr(y_{pse,t}\mid \bm{h}_t,a_t,\overline{\bm{R}}_t = \bm{1}_t,R_{pse,t}=0)}{\pr(y_{pse,t}\mid \bm{h}_t,a_t,\overline{\bm{R}}_t = \bm{1}_t,R_{pse,t}=1)}&= \frac{\exp\{s'_t(\bm{h}_t,a_t)+\gamma'_t y_{pse,t}\}}{E[\exp\{s'_t(\bm{h}_t,a_t)+\gamma'_t y_{pse,t}\} \mid \bm{h}_t,a_t,\overline{\bm{R}}_t = \bm{1}_t, R_{pse,t}=1]}\\
    &= \frac{\exp(\gamma'_t y_{pse,t})}{E\{\exp(\gamma'_t y_{pse,t}) \mid \bm{h}_t,a_t,\overline{\bm{R}}_t = \bm{1}_t, R_{pse,t}=1\}}.
\end{align*}
This proves the equation \eqref{pseudo-obs-mis} in Lemma 1.

\section{Web Appendix G: The $m$-out-of-$n$ bootstrap method for parameter inference in non-regularity scenarios }
The asymptotic properties for stage-1 parameters in Theorem 2 and Theorem 3 in Web Appendix B are based on regularity conditions (i.e., $p_{n-regu} = P\{q_{2,1}(\bm{h}_2;\bm{\psi}_2)=0\} = 0$). However, the inference for the parameters in Q-learning for optimal DTRs may face the challenge of non-regularity \citep{chakraborty2010inference},
which can undermine the validity of the asymptotic distribution provided in Web Appendix B and the ordinary $n$-out-of-$n$ bootstrap method \citep{andrews2000inconsistency}. 
%which can lead to non-smooth estimators and thus inconsistency of the asymptotic distribution provided in Web Appendix B and the ordinary $n$-out-of-$n$ bootstrap method \citep{andrews2000inconsistency}.
Specifically,  in the two-stage scenario, the stage-1 pseudo-outcome and estimators are  non-smooth functions of $\widehat{\bm{\psi}}_2$ when $P\{q_{2,1}(\bm{h}_2;\widehat{\bm{\psi}}_2)=0\} > 0$. As a result, the asymptotic distribution of $(\widehat{\bm{\theta}}_1 - \bm{\theta}_1)$ would be non-normal. Small perturbations in the marginal distribution of $\bm{H}_2$, the parameter $\bm{\psi}_2$, or both, can lead to abrupt differences in the limiting distribution of $\widehat{\bm{\theta}}_1$, and thus make them non-regular. 
In this case, we  followed \cite{chakraborty2013inference} and employed the $m$-out-of-$n$ bootstrap method, which is consistent in non-smooth scenarios,  to construct confidence intervals.

The $m$-out-of-$n$ bootstrap uses a resample size $m$ smaller than the original sample size $n$ to asymptotically avoid the non-negligible probability concentrated at the discontinuous points of $\widehat{Y}_{pse,1}$ and $\widehat{\bm{\theta}}_1$ as functions of $\widehat{\bm{\psi}}_2$ \citep{xu2015regularized}. Specifically, we would draw $B$  bootstrap samples, each with a sample size $m$, and obtain $B$ parameter estimates. The 95\% confidence interval \iffalse can be\fi is then constructed using the 2.5th and 97.5th percentiles of the $B$ parameter estimates. 

%\tbrown{We did not say we use the percentile bootstrap in this section, shall we mention it somewhere? Afterall, the following double bootstrap method to select $m$ used percentile bootstrap.} I don't think we need to emphasize it here. We said we applied $m$-out-of-$n$ bootstrap below, so no need to show the distinction. \tbrown{OK, how about talk more about the $m$-out-of-$n$ and introduce the percentile way implicitly? I have write a draft above} OK


For the choice of resample size $m$, \cite{chakraborty2013inference} proposed a data adaptive method that is directly connected with a measure of non-regularity. 
Specifically, they estimated $p_{n-{regu}}$ by $\widehat{p}_{n-regu} = \widehat{E}(\mathbb{I}[n \{q_{2,1}(\bm{h}_2;\widehat{\bm{\psi}}_2)\}^2 \leq \bm{h}_2\widehat{\Sigma}_{21}\bm{h}_2^{\mathrm{T}} \chi^2_{1,1-\nu}])$, where $\widehat{\Sigma}_{21}$ was the plug-in estimator of $n\text{Cov}(\widehat{\bm{\psi}}_2,\widehat{\bm{\psi}}_2)$, $\chi^2_{1,1-\nu}$ was the $(1-\nu)\times 100$ percentile of the $\chi^2$ distribution with 1 degree of freedom, and $\nu$ was set as $0.001$ in their simulation. We used the same value of $\nu$ in our simulations and MIMIC-III data analysis. The resample size $m$ was set as  $m = \left\lceil n^{\frac{1+\alpha(1-\widehat{p}_{n-regu})}{1+\alpha}} \right\rceil$, where $\lceil x \rceil$ denotes the smallest integer $\geq x$ and $\alpha$ was selected from a grid of candidate values from 0 to 1 using a double bootstrap method.   
Following \cite{chakraborty2013inference}, we applied the following double bootstrap method for choosing $\alpha$ from a grid of candidate values in our simulation and the MIMIC-III data example.


\begin{enumerate}
    \item Denote the estimate of the parameter of interest from original data as $\bm{c}^{\mathrm{T}} \widehat{\bm{\psi}}_1$, where $\bm{c}$ is an arbitrary vector of the same dimension as $\bm{\psi}_1$.
    \item Draw $B_1$ $n$-out-of-$n$ first-level bootstrap samples from original  data and calculate the $\widehat{p}_{n-regu}^{(b_1)}$ with $\nu=0.001$ using the sample where $\overline{\bm{R}}_2^{(b_1)} = \bm{1}_2^{(b_1)}$ (i.e., the sample with complete covariates in the $b_1$th first-level bootstrap \iffalse sample\fi resampling)  %\iffalse i.e.,  the CC sample in the $b_1$th first-stage bootstrap dataset \fi
    for $b_1 = 1,\ldots,B_1$. Fix $\alpha$ at the smallest value in the grid.
    \item Compute the corresponding values of $\widehat{m}^{(b_1)}$ by $$
    \widehat{m}^{(b_1)} = \left\lceil n^{\frac{1+\alpha\left(1-\widehat{p}_{n-regu}^{(b_1)}\right)}{1+\alpha}} \right\rceil,
    $$
    %% where $\lceil x \rceil$ denotes the smallest integer $\geq x$ %\tbrown{(I am wondering whether to add the ceiling function sign in line code 483 when introducing $m$. It is somehow strange that in \cite{chakraborty2013inference}, they did not use the ceiling function in section 2.3.1 when introducing their method, while used the ceiling function 2.3.3 when introducing others' method)} 
     $b_1 = 1,\ldots,B_1$. 
    \item Conditional on each first-level bootstrap sample, draw $B_2$ $m^{(b_1)}$-out-of-$n$ second-level (nested) bootstrap samples and estimate the parameter of interest $\bm{c}^{\mathrm{T}} \widehat{\bm{\psi}}^{(b_1b_2)}_{1,\widehat{m}^{(b_1)}}$, $b_1 = 1,\ldots,B_1$, $b_2 = 1,\ldots,B_2$.
    \item For $b_1=1,...,B_1$, compute the 2.5th and 97.5th percentiles of $\bm{c}^{\mathrm{T}} \widehat{\bm{\psi}}^{(b_1b_2)}_{1,\widehat{m}^{(b_1)}}$; say $\widehat{L}_{db}^{(b_1)}$ and $\widehat{U}_{db}^{(b_1)}$, respectively. %\tbrown{We used the notation $\bm{U}_t$ as '$\bm{U}_t = (\bm{H}_t,A_t)$' and $\bm{l}_t(\bm{z}_t)$ as '$\bm{l}_t(\bm{z}_t) = (l_{t,1}(\bm{z}_t), \ldots, l_{t, L}(\bm{z}_t))$' in the manuscript. Do we need to change the notation $\widehat{L}_{db}^{(b_1)}$ and $\widehat{U}_{db}^{(b_1)}$ here?  } 
    \item Estimate the coverage rate of the bootstrap confidence interval from all the first-level bootstrap datasets as
    $$
    \frac{1}{B_1} \sum_{b_1=1}^{B_1} \mathbb{I} \left\{ 
    \widehat{L}_{db}^{(b_1)} \leq \bm{c}^{\mathrm{T}} \widehat{\bm{\psi}}_1 \leq \widehat{U}_{db}^{(b_1)}
    \right\}.
    $$
    \item  If the above coverage rate is at or above the nominal
    level, then pick the current value of $\alpha$ as the final value. Otherwise, increment $\alpha$ to the next highest value in the grid.
    \item Repeat steps (3)–(7), until the coverage rate of the double bootstrap confidence interval attains the nominal coverage rate or the grid is exhausted.
\end{enumerate}
%\tb{Note that in regular settings, the adaptively selected $m$ would be equal to $n$ and the $m$-out-of-$n$ procedure reduces to the $n$-out-of-$n$ bootstrap.}

As discussed in \cite{chakraborty2013inference}, the number of bootstrap replications in the above double bootstrap procedure to select $\alpha$ (i.e., $B_1$ and $B_2$) does not  need to be equal to the number of  bootstrap replications used for inference (i.e., $B$) after $\alpha$ is selected. Since the kernel regression for estimating the missingness probability models  in the WQ-EE and WQ-SA methods is computationally expensive, we set $B_1=B_2=200$ in our simulation and the MIMIC-III data example.  

\section{Web Appendix H: Simulation}

We conducted four simulation studies to evaluate the finite-sample performance of the WQ-EE and the WQ-SA  methods. In Simulation 1, we focused on the scenario where Assumptions 1-5 were satisfied and demonstrated the performance of the WQ-EE method. In Simulation 2, we evaluated the robustness of the WQ-EE method against the violations of Assumption 5 when (1) the chosen instrumental variable was weakly associated with $R_{pse,t}$ and (2)  there were interactions between $\bm{u}_t$  and $y_{pse,t}$ in the true missingness probability model.  In Simulation 3, we evaluated the proposed simulation-based approach to specifying a plausible range of the sensitivity parameter and the performance of the WQ-SA method based on the specified
range. In Simulation 4, we investigated the influence of an increased number of total stages when covariates are nonignorably missing.  

For comparison, we also assessed three other methods for handling partially missing covariate data in Q-learning: a naive method which ignored partially observed covariates (`naive'), complete case analysis (`CC'), and multiple imputation (`MI') \citep{Little_Rubin_2014}. Q-learning when all the covariates were fully observed (`All') was used as the benchmark. In the naive method, we omitted all terms containing the partially observed covariates in the specified Q-functions. Consequently, the Q-functions were misspecified in this scenario as Assumption 2 was violated. In the CC method, we excluded all patients whose covariates were subject to missingness in estimation. In the MI method, we first imputed missing covariates using predictive mean matching multiple times, which was implemented through the R package \texttt{MICE}. Subsequently, we applied the standard Q-learning algorithm to the imputed data sets. Finally, we calculated the average of the estimated parameters from the imputed datasets to obtain the final estimates of the Q-function parameters. 

All methods were evaluated in two aspects. Firstly, we compared the performances of the methods for estimating the blip function parameters $\bm{\psi}_t$. Secondly, we compared their abilities to identify the true optimal dynamic treatment regime (DTR) and the values of the estimated optimal DTRs. Besides, we also examined the coverage rates of the bootstrap confidence intervals for blip function parameters when the models in the WQ-EE and the WQ-SA methods were correctly specified in Simulation 1 and Simulation 3, respectively.

\subsection{Simulation 1: comparative performance of the WQ-EE method} when Assumptions 1-5 were satisfied
We simulated data from an observational study with two stages of intervention. There were two covariates $X_{t,1}$ and $X_{t,2}$ at stage $t=1,2$. Among these covariates, only $X_{1,2}$ and $X_{2,2}$ were partially observed, and their missingness indicators were $R_1$ and $R_2$, respectively.  Denote $\text{expit}(b) = \exp(b)/\{1+\exp(b)\}$ for $b$ in $\mathcal{R}$. The variables were generated as follows:
\begin{align*}
    & \left( \begin{matrix}X_{1,1}\\X_{2,1}\end{matrix}\right)  \sim N\left(\left( \begin{matrix}0\\0\end{matrix}\right), \left(  \begin{matrix}1 & 0.5\\0.5 & 1 \end{matrix}\right) \right) ;\\
    & X_{1,2} \sim \text{Uniform}(0,2); R_1  \sim \text{Bernoulli} (\{1+\exp(-3+X_{1,2})\}^{-1});\\
    & A_1  \sim 2\times \text{Bernoulli}(\text{expit}(-1+X_{1,1}+X_{1,2}-R_1)) - 1;\\
    & Y_1  =  A_1(\psi_{10}+2+\psi_{11}X_{1,2})+\beta_{10}+\beta_{11} X_{1,1} + \beta_{12} X_{1,2} + \epsilon_1 ; \epsilon_1 \sim N(0,3);\\
    & X_{2,2} \sim \text{Uniform}(0,2);  R_2 \sim  \text{Bernoulli} (\{1+\exp(-1+0.5X_{1,2}-Y_{1}-X_{2,1}-0.5X_{2,2})\}^{-1}) ; \\
    & A_2\sim 2\times \text{Bernoulli} (\text{expit}(-1-X_{1,1}-X_{1,2}+Y_1+X_{2,1}-R_2) ) - 1;\\
    & Y_2 = A_2(\psi_{20}+\psi_{2A}A_1+\psi_{22}X_{2,2}) + \beta_{2A}A_1 + \beta_{22}X_{2,1}+ \beta_{23}X_{2,2} + \epsilon_2; \epsilon_2 \sim N(0,1). 
\end{align*}
%where $ \epsilon_1 \sim N(0,3)$, $\epsilon_2 \sim N(0,1)$.
The true values of the parameters $\bm{\theta}_1 = (\psi_{10},\psi_{11},\beta_{10},\beta_{11},\beta_{12})$ and $\bm{\theta}_2=(\psi_{20},\psi_{2A},\psi_{22},\beta_{2A},\beta_{22},\beta_{23})$ were  $(1,-1,1.5,0.5,-0.5)$ and $(1,-1,1,-1,1,-0.5)$, respectively. We generated $1000$  data sets with sample sizes of $500$ and $2000$.

We considered $Y = Y_1 + Y_2$ as the final outcome. Under this setting, the true optimal treatment rule at stage 2 was $ 2\mathbb{I}(\psi_{20}+\psi_{2A}a_1+\psi_{22}x_{2,2}\geq 0) - 1$. Because $\bm{\theta}_2 = (1,-1,1,-1,1,-0.5)$ and $1-a_1+x_{2,2}\geq 0$, the optimal choice of $A_2$ was always 1. Then, the pseudo-outcome at stage $1$, $Y_{pse,1}$, was equal to $ 1-2A_1+Y_1+X_{2,1}+0.5X_{2,2} $. Therefore, $Q_1(x_{1,1},x_{1,2},a_1) = 1.5 + \beta_{10} + a_1(\psi_{10}+\psi_{11}x_{1,2}) + (0.5+\beta_{11})x_{1,1} - \beta_{12}x_{1,2}$. And the optimal treatment rule at stage 1 was $2\mathbb{I}(\psi_{10}+\psi_{11}X_{1,2}\geq 0) - 1$. Note that the missingness indicators of $X_{1,2}$ and $X_{2,2}$ were directly associated with their own values, thus they were nonignorable missing. $R_1$ only depended on $X_{1,2}$, and $R_2$ only depended on $(X_{1,2},Y_1,X_{2,1},X_{2,2})$,
so the future-independent missingness assumption held. The missingness indicator for $Y_{pse,1}$ was the same as that for $X_{2,2}$, since $Y_{pse,1}=1-2A_1+Y_1+X_{2,1}+0.5X_{2,2}$ and $(A_1, Y_1, X_{2,1})$ were fully observed. Besides, since $P(R_{2}=1  \mid \bm{H}_2)=\{1+\exp(-1+0.5X_{1,2}-Y_{1}-X_{2,1}-0.5X_{2,2})\}^{-1}$, it follows that  $P(R_{pse,1}=1\mid \bm{H}_1,A_1,Y_{pse,1})=\left\{1+\exp(0.5X_{1,2}-2A_1-Y_{pse,1})\right\}^{-1}$, which indicates that $Y_{pse,1}$ was nonignorable missing and $X_{1,1}$ was conditional independent of $R_{pse,1}$ given $(X_{1,2},A_1,Y_{pse,1})$. Since $ X_{1,1}$ was also associated with $Y_{pse,1}$ after conditioning on $(X_{1,2}, A_1)$, we employed $X_{1,1}$ as a nonresponse instrumental variable.

\begin{figure}[htbp]
    \centering
    \includegraphics[scale = 0.5]{1204R2Apsi10.jpg}
    \includegraphics[scale = 0.5]{1204R2Apsi11.jpg}\\
    \includegraphics[scale = 0.5]{1204R2Apsi20.jpg}\includegraphics[scale = 0.5]{1204R2Apsi21.jpg}\\
    \includegraphics[scale = 0.5]{1204R2Apsi22.jpg}
    \caption{Boxplots of blip function parameter estimates from the compared methods in Simulation 1. In each boxplot, white and gray boxes are the results for the sample sizes 500 and 2,000, respectively. The horizontal line marks the true values of the blip function parameter estimates. Note that the naive method did not generate estimates of $\psi_{11}$ and $\psi_{22}$ as $\psi_{11}$ and $\psi_{22}$  are the coefficients of  partially missing $X_{1,2}$  and $X_{2,2}$ in the true treatment rules.  }
    \label{simufig1}
\end{figure}

Web Figure \ref{simufig1} shows the empirical distributions of the estimates of $\bm{\psi}_t= (\psi_{10},\psi_{11}, \psi_{20},\psi_{2A},\psi_{22})$, which directly influence the estimated optimal DTR. At stage 2, the proposed WQ-EE method and the CC method yielded unbiased estimates of the blip function parameters since the future-independent missingness assumption held and the outcome at stage 2 was fully observed. In contrast, the naive and MI methods led to non-negligible biases of the blip function parameters. At stage 1, the estimators in the WQ-EE method remained unbiased while the estimators in the naive, CC, and MI methods displayed relatively large biases, which were not mitigated as the sample size increased. Note that the naive estimates of $\psi_{11}$ and $\psi_{22}$ were unavailable, as we excluded the terms $A_1X_{1,2}$ and $A_2X_{2,2}$ from the specified Q-functions when implementing this method. And the misspecified Q-functions led to the relatively poor performance of the naive estimates. As for the bias of the CC estimates, 
%it can be attributed to differences between the stage 1 Q-function in the complete-case samples and that in the full sample when the pseudo-outcomes are nonignorable missing. That is, obtaining consistent estimates for $\bm{\theta}_1$ directly from the complete-case sample is hindered by the fact that $E(Y_{pse,1}\mid \bm{h}_1,a_1) \neq E(Y_{pse,1}\mid \bm{h}_1,a_1,\overline{\bm{R}}_2 = \bm{1}_2)$. 
it can be attributed to the violation of the ignorable missingness assumption of the pseudo-outcomes, which is necessary for the equation $E(Y_{pse,1}\mid \bm{h}_1,a_1) = E(Y_{pse,1}\mid \bm{h}_1,a_1,\overline{\bm{R}}_2 = \bm{1}_2)$ to hold. 
The bias of the MI estimates stemmed from the nonignorable missingness. It has been discussed that the standard MI method would be biased when data are nonignorably missing \citep{schafer2002missing}.  

The proportion of patients whose optimal treatments in both stages were correctly identified was high when employing the WQ-EE method, as demonstrated in Web Table \ref{simutable1}. With the smaller sample size ($n=500$), the WQ-EE method achieved an average correct classification rate of $89.9\%$ across the 1000 simulated datasets, with a mean final outcome of $2.98$. In general, larger sample sizes led to higher proportions of patients with their true optimal DTR correctly identified. In contrast, the DTRs obtained from the CC and MI methods had much lower correct classification rates and the mean final outcomes under the estimated DTRs were lower than that from the WQ-EE method. 


\begin{table}[htbp]
  \centering
  \caption{Performance of estimated DTRs in Simulation 1}
  \begin{threeparttable}
    \begin{tabular}{lcccc}
    \hline
          & \multicolumn{2}{c}{n=500}     & \multicolumn{2}{c}{n=2000} \\
          \cmidrule(lr){2-3}\cmidrule(lr){4-5}
          & Value & Opt\% & Value & Opt\% \\
          \hline
    All  & 2.988 (0.057) & 0.928 (0.052) & 2.997 (0.028) & 0.964 (0.024) \\
    naive & 2.501 (0.053) & 0.500 (0.022) & 2.500 (0.027) & 0.499 (0.011) \\
    CC    & 2.879 (0.128) & 0.773 (0.111) & 2.898 (0.059) & 0.777 (0.057)  \\
    MI    & 2.832 (0.139) & 0.729 (0.108) & 2.849 (0.071) & 0.730 (0.058) \\
    WQ-EE & 2.974 (0.069) & 0.899 (0.073) & 2.994 (0.029) & 0.949 (0.036) \\
    \hline
    \end{tabular}%
\begin{tablenotes}
    \small
    \item Note: Value, the value of the estimated regime; Opt\%, the percentage of patients correctly classified to their true optimal treatments in both stages; \iffalse Std, the sample standard error. \fi  The numbers outside parentheses represent the sample means, while the numbers inside parentheses represent the sample standard errors.
\end{tablenotes}
  \end{threeparttable}
  \label{simutable1}%
\end{table}%


%We further examined  confidence intervals  of the  blip function parameters constructed by nonparametric bootstrap. 
%We further examined the coverage probability of the constructed confidence intervals. Following the recommendations and numerical studies of  \cite{Shao_Wang_2016}, we constructed the confidence intervals with the bootstrap method due to the complicated form of the asymptotic covariance matrix. Specifically, for each of the 1000 simulated data sets, we obtained 1000 bootstrap samples, and used the 2.5\% and 97.5\% percentiles of the parameter estimates from the bootstrap samples to construct the 95\% confidence intervals for the parameters in this simulated data set. And the coverage probability of the bootstrap confidence intervals was calculated by the percentage of the confidence intervals that covered the true value of corresponding parameter. Web Table \ref{WQ-EE_CI_coverage} shows the results for the coverage probabilities of bootstrap confidence intervals when employing the WQ-EE method in Simulation 1. We see that these coverage probabilities were close to the nominal level $95\%$.
%, indicating that the confidence intervals estimated by the bootstrap method are valid. 

% \tbrown{The general description of the $m$-out-of-$n$ bootstrap takes 2 pages... maybe we should put it as a new Web Appendix section? Or we may abbreviate this part, but the notation $B_1$, $B_2$, and $\nu$ are needed. } OK. Maybe put it in the second web appendix before the proofs or just before simulations? \tbrown{OK, I have moved it just before simulation}

% \tb{We further examined the coverage rates of the constructed confidence intervals. Following the recommendations and numerical studies of  \cite{Shao_Wang_2016}, we constructed the confidence intervals with the bootstrap method due to the complicated form of the asymptotic covariance matrix. As pointed out by a reviewer, the inference for the parameters in Q-learning for optimal DTRs faces the challenge of non-regularity \citep{chakraborty2010inference}, which can lead to non-smooth estimators and thus inconsistency of the ordinary $n$-out-of-$n$ bootstrap method \citep{andrews2000inconsistency}. Specifically,  in the two-stage scenario, the stage-1 pseudo-outcome and estimators are  non-smooth functions of $\widehat{\bm{\psi}}_2$ when $p=P\{q_{2,1}(\bm{h}_2;\bm{\psi}_2)=0\} > 0$. As a result, the asymptotic distribution of $(\widehat{\bm{\theta}}_1 - \bm{\theta}_1)$ would be non-normal and small perturbations in the marginal distribution of $\bm{H}_2$, the parameter $\bm{\psi}_2$, or both, can lead to abrupt differences in the limiting distribution of $\hat{\bm{\theta}}_1$, and thus make them non-regular. For last-stage  parameters, non-regularity is not an issue, allowing the use of the ordinary $n$-out-of-$n$ bootstrap method to construct confidence intervals. 
% %However, for stage-1 parameters, non-regularity may arise when \iffalse $p=P\{q_{2,1}(\bm{h}_2;\hat{\bm{\psi}}_2)=0\} \neq 0$\fi the stage-1 parameter estimators are nonsmooth functions of $\widehat{\bm{\psi}}_2$, 
% %However, non-regularity may arise when the stage-1 parameter estimators are nonsmooth functions of $\widehat{\bm{\psi}}_2$, in which case we  followed \cite{chakraborty2013inference} and employed the $m$-out-of-$n$ bootstrap method. This method uses a bootstrap sample size $m$ smaller than $n$ to asymptotically avoid the non-negligible probability mass concentrated at discontinuous points \citep{xu2015regularized}.
% For stage-1 parameters, we  followed \cite{chakraborty2013inference} and employed the $m$-out-of-$n$ bootstrap method, which is consistent in non-smooth scenarios,  to construct confidence intervals.}


% \tb{For the choice of resample size $m$, \cite{chakraborty2013inference} proposed a data adaptive method that is directly connected with a measure of non-regularity. 
% Specifically, they estimated $p$ by $\widehat{p}_{n-regu} = \widehat{E}(\mathbb{I}[n \{q_{2,1}(\bm{h}_2;\widehat{\bm{\psi}}_2)\}^2 \leq \bm{h}_2\widehat{\Sigma}_{21}\bm{h}_2^{\mathrm{T}} \chi^2_{1,1-\nu}])$, where $\widehat{\Sigma}_{21}$ was the plug-in estimator of $n\text{Cov}(\widehat{\bm{\psi}}_2,\widehat{\bm{\psi}}_2)$, $\chi^2_{1,1-\nu}$ was the $(1-\nu)\times 100$ percentile of the $\chi^2$ distribution with 1 degree of freedom, and $\nu$ was set as $0.001$ in their simulation. The resample size $m$ was set as $m = n^{\frac{1+\alpha(1-\widehat{p}_{n-regu})}{1+\alpha}}$, where $\alpha$ was selected from a grid of candidate values from 0 to 1 using a double bootstrap method.   
% Following \cite{chakraborty2013inference}, we applied the following double bootstrap method for choosing $\alpha$ from a grid of candidate values in our simulation and the MIMIC-III data example.}

% \tb{
% \begin{enumerate}
%     \item Denote the estimate of the parameter of interest from original data as $\bm{c}^{\mathrm{T}} \widehat{\bm{\psi}}_1$, where $\bm{c}$ is an arbitrary vector of the same dimension as $\bm{\psi}_1$.
%     \item Draw $B_1$ $n$-out-of-$n$ first-level bootstrap samples from original  data and calculate the $\widehat{p}_{n-regu}^{(b_1)}$ with $\nu=0.001$ using the sample where $\overline{\bm{R}}_2^{(b_1)} = \bm{1}_2^{(b_1)}$ ( i.e. the sample with complete covariates in the $b_1$th first-level bootstrap sample  \iffalse i.e.,  the CC sample in the $b_1$th first-stage bootstrap dataset \fi), $b_1 = 1,\ldots,B_1$. Fix $\alpha$ at the smallest value in the grid.
%     \item Compute the corresponding values of $\widehat{m}^{(b_1)}$ by $$
%     \widehat{m}^{(b_1)} = n^{\frac{1+\alpha\left(1-\widehat{p}_{n-regu}^{(b_1)}\right)}{1+\alpha}},
%     $$
%     $b_1 = 1,\ldots,B_1$. 
%     \item Conditional on each first-level bootstrap sample, draw $B_2$ $m^{(b_1)}$-out-of-$n$ second-level (nested) bootstrap samples and estimate the parameter of interest $\bm{c}^{\mathrm{T}} \widehat{\bm{\psi}}^{(b_1b_2)}_{1,\widehat{m}^{(b_1)}}$, $b_1 = 1,\ldots,B_1$, $b_2 = 1,\ldots,B_2$.
%     \item For $b_1=1,...,B_1$, compute the 2.5 and 97.5 percentiles of $\bm{c}^{\mathrm{T}} \widehat{\bm{\psi}}^{(b_1b_2)}_{1,\widehat{m}^{(b_1)}}$; say $\widehat{L}_{db}^{(b_1)}$ and $\widehat{U}_{db}^{(b_1)}$, respectively. %\tbrown{We used the notation $\bm{U}_t$ as '$\bm{U}_t = (\bm{H}_t,A_t)$' and $\bm{l}_t(\bm{z}_t)$ as '$\bm{l}_t(\bm{z}_t) = (l_{t,1}(\bm{z}_t), \ldots, l_{t, L}(\bm{z}_t))$' in the manuscript. Do we need to change the notation $\widehat{L}_{db}^{(b_1)}$ and $\widehat{U}_{db}^{(b_1)}$ here?  } 
%     \item Estimate the coverage rate of the bootstrap confidence interval from all the first-level bootstrap datasets as
%     $$
%     \frac{1}{B_1} \sum_{b_1=1}^{B_1} \mathbb{I} \left\{ 
%     \widehat{L}_{db}^{(b_1)} \leq \bm{c}^{\mathrm{T}} \widehat{\bm{\psi}}_1 \leq \widehat{U}_{db}^{(b_1)}
%     \right\}.
%     $$
%     \item  If the above coverage rate is at or above the nominal
%     level, then pick the current value of $\alpha$ as the final value. Otherwise, increment $\alpha$ to the next highest value in the grid.
%     \item Repeat steps (3)–(7), until the coverage rate of the double bootstrap confidence interval attains the nominal coverage rate or the grid is exhausted.
% \end{enumerate}}
% %\tb{Note that in regular settings, the adaptively selected $m$ would be equal to $n$ and the $m$-out-of-$n$ procedure reduces to the $n$-out-of-$n$ bootstrap.}


We further examined the coverage rates of the constructed confidence intervals. Following the recommendations and numerical studies of  \cite{Shao_Wang_2016}, we constructed the confidence intervals with the bootstrap method due to the complicated form of the asymptotic covariance matrix. For \iffalse last-stage\fi stage-2  parameters, non-regularity is not an issue, allowing the use of the ordinary $n$-out-of-$n$ bootstrap method to construct confidence intervals. For stage-1 parameters, we employed the $m$-out-of-$n$ bootstrap method described in Web Appendix G to construct confidence intervals.  %In our simulations, 
%When determining the tuning parameter $\alpha$  for  the bootstrap resampling size $m$ for each simulated dataset,


Because the kernel regression for estimating missingness probabilities in the WQ-EE method slows down the computation, the adaptive selection of $\alpha$ for determining the bootstrap resample size $m$ in each simulated dataset is computationally prohibitive. To reduce the computation burden in our simulations, we set $B_1 = B_2=200$ in the double bootstrap method and set $\alpha$ at the fixed value that was determined by calculating the average of the $\alpha$ values adaptively selected \iffalse for\fi from 10 simulated datasets, each with a sample size 500. The parameter of interest in the double bootstrap procedure was set as $\psi_{10}$. 

With the determined $\alpha$, we conducted 500 bootstrap resampling for each of the 1000 simulated datasets to calculate the $n$-out-of-$n$ and $m$-out-of-$n$ percentile bootstrap confidence intervals for the stage-2 and stage-1 blip function parameters, respectively. And the coverage rate of the bootstrap confidence intervals was calculated by the percentage of the confidence intervals that covered the true value of corresponding parameter.

Web Table \ref{WQ-EE_CI_coverage} shows the results for the coverage rates of the bootstrap confidence intervals when employing the WQ-EE method in Simulation 1. We see that these coverage rates were close to the nominal level $95\%$. 



% \begin{table}[htbp]
%   \centering
%   \caption{Coverage probability for the bootstrap 95\% confidence intervals when employing the WQ-EE method in Simulation 1}
%   \begin{threeparttable}
%     \begin{tabular}{lccccc}
%     \hline
%      & $\psi_{10}$ & $\psi_{11}$ & $\psi_{20}$ & $\psi_{2A}$ & $\psi_{22}$ \\
%     \hline
%      Sample size 500 & 0.923 & 0.954 & 0.925 & 0.971 & 0.938 \\
%      Sample size 2000 & 0.939 & 0.945 & 0.940 & 0.956 & 0.944 \\
%     \hline
%     \end{tabular}
% \begin{tablenotes}
%     \small
%     \item Note: The 95\% confidence intervals are constructed by bootstrap percentiles.
% \end{tablenotes}
%   \end{threeparttable}
%   \label{WQ-EE_CI_coverage}%
% \end{table}%

\begin{table}[htbp]
  \centering
  \caption{Coverage rates of  the 95\% bootstrap  confidence intervals when employing the WQ-EE method in Simulation 1}
  \begin{threeparttable}
    \begin{tabular}{lcccccc}
    \hline
    &$m$ & $\psi_{10}$ & $\psi_{11}$ & $\psi_{20}$ & $\psi_{2A}$ & $\psi_{22}$ \\
    \hline
     Sample size 500 & 473 & 0.947 & 0.937 & 0.951 & 0.941 & 0.946 \\
     Sample size 2000 & 1956 & 0.947 & 0.945 & 0.953 & 0.950 & 0.945 \\
    \hline
    \end{tabular}
\begin{tablenotes}
    \small
    \item Note: $m$, the average  resample size in the $m$-out-of-$n$ bootstrap method in 1000 simulated datasets, rounded to integer.   For stage-1 parameters ($\psi_{10}, \psi_{11}$), the 95\% confidence intervals \iffalse are\fi were constructed by $m$-out-of-$n$ bootstrap percentiles. For stage-2 parameters ($\psi_{20}, \psi_{2A}, \psi_{22}$), the 95\% confidence intervals \iffalse are\fi were constructed by standard ($n$-out-of-$n$) bootstrap percentiles. 
\end{tablenotes}
  \end{threeparttable}
  \label{WQ-EE_CI_coverage}%
\end{table}%


      


\subsection{Simulation 2: comparative performance of the WQ-EE method when Assumption 5 is violated}

To investigate the robustness of the proposed WQ-EE method when Assumption 5 was violated, we kept all other data generating mechanisms unchanged, while altering the true model for $R_{pse,1}$ to $P(R_{pse,1}=1) = \left\{1+\exp(\gamma_z X_{1,1} + 0.5X_{1,2}-2A_1-Y_{pse,1})\right\}^{-1}$. When $\gamma_z=0$, the data generating mechanism was the same as that in Simulation 1 and Assumption 5 was satisfied. However, when $\gamma_z \neq 0$, $X_{1,1}$ was associated  with $R_{pse,1}$ conditional on $(X_{1,2},A_1,Y_{pse,1})$. Thus there were no nonresponse instrumental variables. We set $\gamma_z \in \{-0.4,-0.2,0.2,0.4\}$ such that $X_{1,1}$ was weakly associated with $P(R_{pse,1}=1)$ conditional  on $X_{1,2}$ and $Y_{pse,1}$.  We generated 1000 Monte Carlo samples with sample size $n=500$. 

Since the future-independent missingness assumption still held in Simulation 2, the estimators for $(\psi_{20},\psi_{2A},\psi_{22})$  in the WQ-EE method were still consistent. We report the performance of the estimators for $(\psi_{10},\psi_{11})$ and the value and correct classification rate of the estimated optimal DTRs. The results were calculated by averaging over  1000 simulated data sets and are shown in Web Table~\ref{simu_IV_robust}. We see that the proposed WQ-EE method was robust when $X_{11}$ was weakly related to $R_{pse,1}$. Despite the declining performance of the corresponding optimal DTR as the absolute value of $\gamma_z$ increased, the WQ-EE method consistently outperformed the other three methods with higher mean final outcomes and correct classification rates.

\begin{table}[htbp] 
\centering
\begin{threeparttable}
  \caption{Performance of the WQ-EE method without valid nonresponse instrumental variable in comparison with alternative methods in Simulation 2}\label{simu_IV_robust}
    \begin{tabular}{lcccccccc}
    \hline
           & \multicolumn{2}{c}{$\psi_{10}$}     & \multicolumn{2}{c}{$\psi_{11}$}  & \multicolumn{2}{c}{Value} &  \multicolumn{2}{c}{Opt\%} \\
          \cmidrule(lr){2-3}\cmidrule(lr){4-5}\cmidrule(lr){6-7}\cmidrule(lr){8-9}
          Method  & Bias & Std & Bias & Std & Mean & Std  & Mean & Std \\
          \hline
          \multicolumn{9}{c}{$\gamma_z = -0.4$} \\
    naive  & -1.005 & 0.168 &  \textemdash  &\textemdash  &2.498 &0.054 &0.499 &0.022 \\
    CC    & 0.068 & 0.188 & -0.447 & 0.277 & 2.881 & 0.131 & 0.773 & 0.113  \\
    MI    & 0.166 & 0.169 & -0.588 & 0.242 & 2.838 & 0.144 & 0.738 & 0.116 \\
    WQ-EE & -0.042 & 0.237 & 0.113 & 0.318 & 2.971 & 0.078 & 0.896 & 0.080 \\
              \multicolumn{9}{c}{$\gamma_z = -0.2$} \\
    naive & -1.008 & 0.166 &  \textemdash &\textemdash&2.498 &0.054 & 0.500 & 0.023 \\
    CC    & 0.062 & 0.184 & -0.443 & 0.266 & 2.885 & 0.126 & 0.781 & 0.113  \\
    MI    & 0.161 & 0.165 & -0.591 & 0.231 & 2.838 & 0.142 & 0.733 & 0.113 \\
    WQ-EE & -0.043 & 0.243 & 0.090 & 0.315 & 2.972 & 0.081 & 0.898 & 0.076 \\
              \multicolumn{9}{c}{$\gamma_z = 0.2$} \\
    naive  & -1.000 & 0.155 &  \textemdash &\textemdash& 2.502 & 0.052 & 0.499 & 0.022 \\
    CC    & 0.073 & 0.185 & -0.482 & 0.253 & 2.871 & 0.125 & 0.766 & 0.110  \\
    MI    & 0.174 & 0.164 & -0.627 & 0.227 & 2.816 & 0.144 & 0.718 & 0.113 \\
    WQ-EE & -0.040 & 0.239 & 0.031 & 0.310 & 2.974 & 0.070 & 0.899 & 0.072 \\
              \multicolumn{9}{c}{$\gamma_z = 0.4$} \\
    naive  & -1.007 &0.151  & \textemdash &\textemdash& 2.501 & 0.053 & 0.499 & 0.023 \\
    CC    & 0.088 & 0.185 & -0.495 & 0.274 & 2.862 & 0.132 & 0.755 & 0.114  \\
    MI    & 0.183 & 0.163 & -0.639 & 0.238 & 2.808 & 0.146 & 0.705 & 0.114 \\
    WQ-EE & -0.020 & 0.233 & 0.053 & 0.321 & 2.968 & 0.074 & 0.895 & 0.078 \\
    \hline
    \end{tabular}%
   \begin{tablenotes}
       \small
       \item Note: Value, the value of the estimated regime; Opt\%, the percentage of patients correctly classified to their true optimal treatments in both stages; Bias, the bias of the sample mean; Std, the sample standard error; Mean, the sample mean.
   \end{tablenotes}
\end{threeparttable}
\end{table}


We also considered scenarios where there were interactions between $y_{pse,t}$ and $\bm{u}_t$ in the true missingness probability model. The true model for $R_{pse,1}$ was modified as $P(R_{pse,1}=1) = \left[1+\exp\{0.5X_{1,2}-2A_1-Y_{pse,1}+\gamma_{uy}(X_{1,2}-1)Y_{pse,1}\}\right]^{-1}$, while the data generating mechanism for other covariates was the same as that in Simulation 1. The specific form for the interaction, $(X_{1,2}-1)Y_{pse,1}$, was chosen to ensure that the missing data proportions were largely similar between Simulations 1 and 2. 
%\JS{Here I use '$(X_{12}-1)Y_{pse,1}$' to make sure the missing proportion will not change too much when $\gamma_{UY}\neq 0$} 
We set $\gamma_{uy} \in \{-0.2,-0.1,0.1,0.2\}$ by  allowing weak interactions between $X_{1,2}$ and $Y_{pse,1}$.
%and the proposed method based on $P(R_{pse,t}=1\mid \bm{u}_t,y_{pse,t},\overline{\bm{R}}_t = \bm{1}_t) = [1+\exp(s_t(\bm{u}_t)+\gamma_t y_{pse,t})]^{-1}$ is exposed to model misspecification in terms of the parametric function for $y_{pse,1}$.  
We again generated  1000 simulated data sets with sample size $n=500$. The results of the compared methods are shown in Web Table~\ref{simu_inter_robust}. 
The WQ-EE method exhibited superior performance compared to the other three methods despite the introduction of interaction terms involving $Y_{pse,1}$, showcasing its robustness in the presence of mild model misspecification.


\begin{table}[htbp]
  \centering
    \caption{Performance of the WQ-EE method when interaction terms between $y_{pse,t}$ and $\bm{u}_t$ were incorrectly omitted, in comparison with alternative methods  in Simulation 2} \label{simu_inter_robust}
\begin{threeparttable}
    \begin{tabular}{lcccccccc}
    \hline
           & \multicolumn{2}{c}{$\psi_{10}$}     & \multicolumn{2}{c}{$\psi_{11}$}  & \multicolumn{2}{c}{Value} &  \multicolumn{2}{c}{Opt\%} \\
          \cmidrule(lr){2-3}\cmidrule(lr){4-5}\cmidrule(lr){6-7}\cmidrule(lr){8-9}
          Method  & Bias & Std & Bias & Std & Mean & Std  & Mean & Std \\
          \hline
          \multicolumn{9}{c}{$\gamma_{uy} = -0.2$} \\
naive & -1.011 & 0.160  & \textemdash &\textemdash& 2.500 & 0.053 & 0.500 & 0.022 \\
CC & 0.124 & 0.184 & -0.514 & 0.269 & 2.863 & 0.139 & 0.760 & 0.115 \\
MI & 0.215 & 0.162 & -0.656 & 0.233 & 2.808 & 0.152 & 0.712 & 0.115 \\
WQ-EE & -0.067 & 0.239 & 0.075 & 0.324 & 2.972 & 0.072 & 0.898 & 0.073 \\
\multicolumn{9}{c}{$\gamma_{uy} = -0.1$} \\
naive & -1.002 & 0.153  & \textemdash &\textemdash& 2.500 & 0.053 & 0.499 & 0.022\\
CC & 0.084 & 0.182 & -0.474 & 0.270 & 2.871 & 0.132 & 0.764 & 0.116 \\
MI & 0.185 & 0.159 & -0.624 & 0.230 & 2.820 & 0.141 & 0.715 & 0.115 \\
WQ-EE & -0.062 & 0.228 & 0.065 & 0.318 & 2.971 & 0.072 & 0.897 & 0.075 \\
\multicolumn{9}{c}{$\gamma_{uy} = 0.1$} \\
naive & -1.002 & 0.157  & \textemdash &\textemdash& 2.499 & 0.053 & 0.500 & 0.022\\
CC & 0.046 & 0.182 & -0.441 & 0.258 & 2.886 & 0.119 & 0.780 & 0.105 \\
MI & 0.147 & 0.163 & -0.595 & 0.231 & 2.832 & 0.143 & 0.734 & 0.107 \\
WQ-EE & -0.021 & 0.233 & 0.022 & 0.307 & 2.974 & 0.074 & 0.898 & 0.077 \\
\multicolumn{9}{c}{$\gamma_{uy} = 0.2$} \\
naive & -1.010 & 0.160  & \textemdash &\textemdash& 2.497 & 0.051 & 0.499 & 0.023 \\
CC & 0.046 & 0.188 & -0.444 & 0.263 & 2.883 & 0.124 & 0.772 & 0.112 \\
MI & 0.146 & 0.168 & -0.590 & 0.234 & 2.835 & 0.140 & 0.732 & 0.111 \\
WQ-EE & 0.023 & 0.242 & -0.011 & 0.308 & 2.972 & 0.079 & 0.895 & 0.082 \\
    \hline
    \end{tabular}%
\begin{tablenotes}
    \small
    \item Note: Value, the value of the estimated regime; Opt\%, the percentage of patients correctly classified to their true optimal treatments in both stages; Bias, the bias of the sample mean; Std, the sample standard error; Mean, the sample mean.
\end{tablenotes}
\end{threeparttable}
\end{table}%


%Overall, the simulation results in stages 1 and 2 demonstrated that the proposed WQ-EE method performed well under correct model specification and mild model misspecification, in comparison with the naive, CC, and MI methods. 

\subsection{Simulation 3: The performance of the WQ-SA method and the proposed simulation-based approach to calibrating the magnitude of the sensitivity parameters}

To evaluate the performance of the WQ-SA method and the proposed simulation-based approach to calibrating the magnitude of the sensitivity parameters, we simulated  data from a 2-stage scenario with two covariates at each stage. At stage $t = 1, 2$, there existed two covariates, $X_{t,1}$ and $X_{t,2}$. Among these covariates, only $X_{1,2}$ and $X_{2,2}$ were partially observed, with their corresponding missingness indicators denoted as $R_1$ and $R_2$, respectively. The data-generation mechanism was as follows:

\begin{align*}
    & \left( \begin{matrix}X_{1,1}\\X_{2,1}\end{matrix}\right)  \sim N\left(\left( \begin{matrix}0\\0\end{matrix}\right), \left(  \begin{matrix}1 & 0.5\\0.5 & 1 \end{matrix}\right) \right) ;\\
    & X_{1,2} \sim \text{Uniform}(0,2); \\
    &R_1  \sim \text{Bernoulli} (\{1+\exp(-1-X_{1,2})\}^{-1});\\
    & A_1  \sim 2\times \text{Bernoulli}(\text{expit}(1-X_{1,1}-X_{1,2}+R_1)) - 1;\\
    & Y_1  =  A_1(\psi_{10}'+2+\psi_{11}'X_{1,2})+\beta_{10}'+\beta_{11}' X_{1,1} + \beta_{12}' X_{1,2} + \epsilon_1' ; \epsilon_1' \sim N(0,3);\\
    & X_{2,2} \sim \text{Uniform}(0,2); \\
    &R_2 \sim  \text{Bernoulli} ([1+\exp\{2.5-\exp(2X_{1,1})-0.5X_{1,2}+A_1+Y_{1}+X_{2,1}+0.5X_{2,2}\}]^{-1}) ; \\
    & A_2\sim 2\times \text{Bernoulli} (\text{expit}(-1+X_{1,1}-2A_1+X_{2,1}) ) - 1;\\
    & Y_2 = A_2(\psi_{20}'+\psi_{2A}'A_1+\psi_{22}'X_{2,2}) + \beta_{2A}'A_1 + \beta_{22}'X_{2,1}+ \beta_{23}'X_{2,2} + \epsilon_2'; \epsilon_2' \sim N(0,1). 
\end{align*}
The true values of the parameters $\bm{\theta}_1' = (\psi_{10}',\psi_{11}',\beta_{10}',\beta_{11}',\beta_{12}')$ and $\bm{\theta}_2' = (\psi_{20}',\psi_{2A}',\beta_{22}',\beta_{2A}',\beta_{22}',\beta_{23}')$ were $(-1,1,-2,-0.5,-1)$ and $(1,-1,1,-1,1,-0.5)$, respectively. The final outcome was set to be $Y=Y_1+Y_2$. 

%\tp{(The DGM for $\bm{X}$ and $Y_2$ is unchanged, while others change. In the the beginning I tried to just change $R_2$, but the results are not satisfying by 1. in SSA, $\gamma=0$ have larger p-values than $\gamma=1$ (usually), 2. in WQ-SA, $\gamma=2$ have better performance than $\gamma=1$ (often, can be resolved by changing the bandwidth in kernel regression), 3. CC and MI perform too well with correct classification rate above 0.8 (sometimes). After trying for 2.5 days, I gave up this way and changed DGM for other covariates.  )}

The true optimal treatment rule at stage 2 was $2\mathbb{I}(1-a_1+x_{2,2}\geq 0)-1$. Since $a_1\in \{-1,1\}$ and $x_{2,2} \geq 0$, the optimal choice of $A_2$ was always 1. With this optimal treatment at stage 2, $Y_{pse,1} = 1-2A_1+Y_1+X_{2,1}+0.5X_{2,2}$. Then, 
%$Q_1(x_{1,1},x_{1,2},a_1) = 1.5+ \beta_{10}+a_1(\psi_{10}+\psi_{11}x_{1,2})+(0.5+\beta_{11})x_{1,1}-\beta_{12}x_{1,2}$ 
the optimal treatment rule at stage 1 was $2\mathbb{I}(\psi_{10}+\psi_{11}x_{1,2}\geq 0)-1$. Since the missingness indicators of $X_{1,2}$ and $X_{2,2}$ were directly related to their own values, they were nonignorably missing. The future-independent missingness assumption held because $R_1$ only depended on $X_{1,2}$, and $R_2$ only depended on $(X_{1,1},X_{1,2},A_1,Y_1,X_{2,1},X_{2,2})$. Because $Y_{pse,1} = 1-2A_1+Y_1+X_{2,1}+0.5X_{2,2}$, $R_{pse,1} = R_2$, $P(R_{pse,1} =1) = \{ 1.5-\exp(2X_{1,1})-0.5X_{1,2}+3A_1+Y_{pse,1} \}^{-1}$, $Y_{pse,1}$ was nonignorably missing and there was no nonresponse instrumental variable. And the true value of the sensitivity parameter was  $\gamma_1'=1$ in this scenario.


We first evaluated the performance of the simulation-based approach to calibrating the magnitude of the sensitivity parameter. 
%We generated 1000 simulated data sets with sample size $n=500$ and $2000$, respectively. 

As discussed in Web Appendix A, we assumed the sign of the sensitivity parameter $\gamma_1'$ was known and chose $[0,7]$ as an initial range for $\gamma_1'$. To zoom in to a more plausible range, we employed the simulation-based approach proposed in Web Appendix A. Specifically, we first generated 1000 data sets with sample size $n=500$ and $2000$, respectively. For each data set, to evaluate the candidate values of $\gamma_1'$, we implemented the steps described in the end of  Web Appendix A. Particularly, in Step (2), $\pr(y_{pse,1}\mid x_{1,1},x_{1,2},a_1,R_1=1,R_2=1)$ was obtained by combing the kernel regression estimates of $E(y_{pse,1}\mid x_{1,1},x_{1,2},a_1,R_1=1,R_2=1)$ and its residual density estimates, which were obtained by kernel density estimation method with the R function \texttt{density}. And $E\{\exp(\gamma_1'y_{pse,1}) \mid x_{1,1},x_{1,2},a_1,R_1=1,R_2=1 \}$ was also estimated by kernel regression methods. In Step (3), based on the estimated $\pr(y_{pse,1}\mid x_{1,1},x_{1,2},a_1,R_1=0,R_2=1)$, we imputed $Y_{pse,1}^{\gamma_1',imp}$ for 1000 Monte Carlo replication times. And in Step (4), we estimated $\pr(y_{pse,1}^{\gamma_1',*}\mid x_{1,1},x_{1,2},a_1,R_1=1)$ with the correctly specified Q-function for the conditional mean and R function \texttt{density} for the residuals. Further, we generated the replications of observed pseudo-outcomes and compared them with the observed pseudo-outcome estimates using the Wilcoxon ranked sum test in Steps (5) and (6), respectively.
Finally, we calculated the medians of the Wilcoxon ranked sum test p-values from the 1000 Monte Carlo replicates for each simulated data set, and we reported the summary statistics of these  p-value medians from the 1000 simulated data sets and the  probability to be included in the  range of plausible values determined by the threshold of 0.05 in Web Table \ref{simu_sen_spec}.    


\begin{table}[htbp]
    \centering
\begin{threeparttable}
    \caption{Performance of the proposed simulation-based approach to determining a plausible range of the sensitivity parameter 
 in Simulation 3}      \label{simu_sen_spec}
    \begin{tabular}{lcccccccc}
    \hline
    $\gamma_t'$ & 0 & 1 & 2 & 3 & 4 & 5 & 6  & 7 \\
    \hline
                  \multicolumn{9}{c}{Sample size = 500} \\
    minimum & 0.211 & 0.216 & 0.049 & 0.006 & 0.002 & 0.000 & 0.000 & 0.000\\
    median &0.591 & 0.615 & 0.452& 0.273& 0.201& 0.143& 0.101& 0.073\\
    mean & 0.578 &0.601 &0.440& 0.297 & 0.232 & 0.185 & 0.150 & 0.122\\
    maximum & 0.754 & 0.756 & 0.744 & 0.762 & 0.695 & 0.665 & 0.671 & 0.662\\
    inclusion probability  & 1.000 &  1.000 &  0.997 &  0.966 &  0.895 &  0.814 &  0.722 &  0.605\\
                   \multicolumn{9}{c}{Sample size = 2000} \\
    minimum & 0.105 & 0.338 & 0.001 & 0.000 & 0.000 & 0.000 & 0.000& 0.000 \\
    median & 0.377 & 0.624 & 0.045 & 0.010 & 0.003 & 0.001 & 0.000 &0.000\\
    mean & 0.381 & 0.618 & 0.061 & 0.015 & 0.008 & 0.005 & 0.003 & 0.002\\
    maximum & 0.681 & 0.757 & 0.342 & 0.225 & 0.089 & 0.108 & 0.094 & 0.072\\
    inclusion probability  & 1.000 & 1.000 & 0.448 & 0.055 & 0.026 & 0.013 & 0.012 & 0.007\\
    \hline
    \end{tabular}
    \begin{tablenotes}
        \small
        \item Note: minimum, minimum of p-value medians in the 1000 data sets; median, median of p-value medians in the 1000 data sets;   
        mean, mean of p-value medians in the 1000 data sets;
        maximum, maximum of p-value medians in the 1000 data sets;
        inclusion probability, the probability to be included in the selected range based on threshold 0.05.
    \end{tablenotes}
\end{threeparttable}
\end{table}

% \tp{How about we do not say the selected range in the simulation but show the WQ-SA for all $\gamma_1'$ in [0,7], we actually select 1000 different ranges for $\gamma_1'$ in the 1000 data sets. Some would be [0,2], some would be [0,7]. I think it may be easier to understand that we use the same way to select the range in simulation and real data. And in sample size 500, we actually included [0,7] based on the inclusion probability. By the way, for sample size 500, we would included $\gamma_1'=12$ based on the inclusion probability }

% OK. I think it is fair to not to link the range determination procedure with the following simulations for the WQ-SA method. I will cut some sentences.

% \tp{We may say the selected range varies across different datasets (Actually the word 'inclusion probability' already indicates that the range is determined for each data set.)   } OK, i clarified that the threshold applied to each data set.
% \tp{OK, thanks, the code for $\gamma_1'=4,5,6,7$ for WQ-SA is running, it should be ok tomorrow}

 We see that for both sample sizes, the summary statistics corresponding to the true value $\gamma_1'=1$ were largest among candidate values of $\gamma_1'$, and they decreased as the specified sensitivity parameter moving away from its true value. In analogy to the conventional significance level $0.05$ for hypothesis testing based on p-values, we considered a threshold of the p-value medians 0.05 to determine a plausible range  of $\gamma_1'$ for each simulated data set, where a value of $\gamma_1'$  was included in the determined range if its corresponding p-value median was greater than 0.05.  The true value for $\gamma_1'$ was $100\%$ included in the determined ranges when applying this threshold to the 1000 data sets. This indicated that applying the 0.05 threshold for the p-value medians can ensure that the determined plausible ranges  would cover the true value of $\gamma_1'$.

With larger sample sizes, the summary statistics of the p-value medians for the true value of $\gamma_1'$ was stable while those for other values of $\gamma_1'$ decreased. With a sample size of 2000, $\gamma_1'=3$ was included in the determined plausible range for sensitivity parameter in 54 out of 1000 data sets, and larger values of $\gamma_1'$ were selected even  less frequently.% Therefore, we used the range  $[0,3]$ in the subsequent investigation on   the performance of the WQ-SA method. }
%\tb{Note that the choice of this range based on inclusion probabilities was different from what can be applied to  real data settings without replication  data. However, the results in Web Table~\ref{simu_sen_spec} demonstrated that the choice of the threshold  0.05 worked well to include the true value of  $\gamma_1'$. }  




We compared the performance of the WQ-SA method with $\gamma_1' \in [0,7]$ with other methods.  First we assessed their performance in estimating blip function parameters and optimal DTRs. Similar to Simulations 1 and 2, the future-independent missingness assumption still held in Simulation 3, and we reported the performance of the estimators for $(\psi_{10}',\psi_{11}')$ and the values and correct classification rates of the estimated optimal DTRs. The results were calculated by averaging over 1000 simulated datasets and are shown in Web Table \ref{WQ-SA_0506}.  


\begin{table}[htbp] 
\centering
\begin{threeparttable}
  \caption{Performance of the WQ-SA method in Simulation 3}\label{WQ-SA_0506}
    \begin{tabular}{lcccccccc}
    \hline
           & \multicolumn{2}{c}{$\psi_{10}'$}     & \multicolumn{2}{c}{$\psi_{11}'$}  & \multicolumn{2}{c}{Value} &  \multicolumn{2}{c}{Opt\%} \\
          \cmidrule(lr){2-3}\cmidrule(lr){4-5}\cmidrule(lr){6-7}\cmidrule(lr){8-9}
          Method  & Bias & Std & Bias & Std & Mean & Std  & Mean & Std \\
          \hline
          \multicolumn{9}{c}{Sample size = 500} \\
All & 0.016 & 0.234 & -0.011 & 0.158 & -1.018 & 0.045 & 0.923 & 0.055\\
naive & 1.171 & 0.162 &  \textemdash & \textemdash & -1.498 & 0.035& 0.503& 0.023 \\
CC & -0.319 & 0.298 & -0.078 & 0.209 & -1.131 & 0.121 & 0.767 & 0.114 \\
MI & -0.269 & 0.264 & -0.178 & 0.186 & -1.184 & 0.138 & 0.721 & 0.113\\
WQ-SA ($\gamma_1'=0$) & -0.037 & 0.213 & -0.395 & 0.298 & -1.134 & 0.124 & 0.760 & 0.111\\
WQ-SA ($\gamma_1'=1$) & -0.056 & 0.241 &  0.034 & 0.335 & -1.033 & 0.071 & 0.893 & 0.084\\
WQ-SA ($\gamma_1'=2$) & -0.049 & 0.251 &  0.083 & 0.347 & -1.034 & 0.071 & 0.891 & 0.084\\
WQ-SA ($\gamma_1'=3$) & -0.056 & 0.252 &  0.062 & 0.339 & -1.036 & 0.070 & 0.889 & 0.086\\
WQ-SA ($\gamma_1'=4$) & -0.072 & 0.256 &  0.058 & 0.344 & -1.037 & 0.077 & 0.887 & 0.088\\
WQ-SA ($\gamma_1'=5$) & -0.070 & 0.264 &  0.025 & 0.349 & -1.041 & 0.077 & 0.883 & 0.091\\
WQ-SA ($\gamma_1'=6$) & -0.053 & 0.263 & -0.015 & 0.361 & -1.042 & 0.079 & 0.880 & 0.093\\
WQ-SA ($\gamma_1'=7$) & -0.060 & 0.271 & -0.023 & 0.359 & -1.044 & 0.080 & 0.879 & 0.095\\
\multicolumn{9}{c}{Sample size = 2000} \\
All & 0.005 & 0.118 & -0.002 & 0.080 & -1.003 & 0.020 & 0.965 & 0.025 \\
naive & 1.173 & 0.077 & \textemdash & \textemdash & -1.500 & 0.014 & 0.500 & 0.011 \\
CC & -0.324 & 0.153 & -0.075 & 0.106 & -1.101 & 0.053 & 0.774 & 0.057 \\
MI & -0.296 & 0.134 & -0.152 & 0.096 & -1.148 & 0.065 & 0.731 & 0.056\\
WQ-SA ($\gamma_1'=0$) & -0.036 & 0.102 & -0.421 & 0.143 & -1.120 & 0.056 & 0.756 & 0.055\\
WQ-SA ($\gamma_1'=1$) & -0.026 & 0.140 &  0.035 & 0.182 & -1.007 & 0.022 & 0.948 & 0.037\\
WQ-SA ($\gamma_1'=2$) & -0.019 & 0.132 & -0.036 & 0.177 & -1.008 & 0.023 & 0.943 & 0.043\\
WQ-SA ($\gamma_1'=3$) & -0.021 & 0.158 & -0.019 & 0.205 & -1.009 & 0.024 & 0.939 & 0.043\\
WQ-SA ($\gamma_1'=4$) & -0.020 & 0.163 & -0.066 & 0.200 & -1.012 & 0.026 & 0.931 & 0.051\\
WQ-SA ($\gamma_1'=5$) & -0.033 & 0.165 & -0.098 & 0.206 & -1.019 & 0.029 & 0.915 & 0.056\\
WQ-SA ($\gamma_1'=6$) & -0.032 & 0.164 & -0.137 & 0.203 & -1.027 & 0.035 & 0.897 & 0.062\\
WQ-SA ($\gamma_1'=7$) & -0.051 & 0.168 & -0.139 & 0.209 & -1.031 & 0.037 & 0.887 & 0.063\\
    \hline
    \end{tabular}%
   \begin{tablenotes}
       \small
       \item Note: WQ-SA ($\cdot$), the WQ-SA method based on corresponding sensitivity parameters; Value, the value of the estimated regime; Opt\%, the percentage of patients correctly classified to their true optimal treatments in both stages; Bias, the bias of the sample mean; Std, the sample standard error; Mean, the sample mean.
   \end{tablenotes}
\end{threeparttable}
\end{table}

We see that when the specified sensitivity parameter was equal to its true value, the WQ-SA method had the best performance. Notably,  when $\gamma_t'=0$ was specified, it was equivalent to assuming ignorable missingness, that is,  the missingness of $Y_{pse,1}$ was conditionally independent with its own values. Except for this specification, the WQ-SA method exhibited relative robustness to the sensitivity parameter specification and consistently outperformed the naive, CC, and MI methods with higher mean final outcomes and correct classification rates under the estimated optimal DTRs.  

Finally, we examined the coverage rates of the confidence intervals for blip function parameters when employing the WQ-SA method with the correctly specified sensitivity parameter value $\gamma_t'=1$. 
%Similar to Simulation 1, the confidence intervals were constructed by bootstrap percentiles.
Similar to Simulation 1, the confidence intervals for stage-2 parameters were constructed by ordinary $n$-out-of-$n$ bootstrap method, while those for stage-1 parameters were constructed by $m$-out-of-$n$ bootstrap method described in Web Appendix G. And the bootstrap resample sizes were selected in the  same way as in Simulation 1. 
%To reduce the computational burden, we also used a fixed $\alpha$, which was selected as the average value of those adaptively chosen from 10 datasets using the double bootstrap method with $B_1 = B_2 = 200$ and a sample size of 500 as performed in Simulation 1, in 1000 Monte Carlo simulated datasets. 
Web Table \ref{WQ-SA_CI_coverage} shows the coverage rates of the bootstrap confidence intervals, which were close to the nominal level $95\%$.   

% \begin{table}[htbp]
%   \centering
%   \caption{Coverage probability for the estimated 95\% confidence intervals when employing the WQ-SA method with correct sensitivity parameter in Simulation 3}
% \begin{threeparttable}
%     \begin{tabular}{lccccc}
%     \hline
%      & $\psi_{10}'$ & $\psi_{11}'$ & $\psi_{20}'$ & $\psi_{2A}'$ & $\psi_{22}'$ \\
%     \hline
%      Sample size 500 & 0.936 & 0.941 & 0.948 & 0.927 & 0.938 \\
%      Sample size 2000 & 0.940 & 0.937 & 0.941 & 0.938 & 0.935 \\
%     \hline
%     \end{tabular}
% \begin{tablenotes}
%     \small
%     \item Note: The 95\% confidence intervals were constructed by bootstrap percentiles.
% \end{tablenotes}
%   \end{threeparttable}
%   \label{WQ-SA_CI_coverage}%
% \end{table}%

\begin{table}[htbp]
  \centering
  \caption{Coverage rates of  the bootstrap 95\% confidence intervals when employing the WQ-SA method in Simulation 3}
  \begin{threeparttable}
    \begin{tabular}{lcccccc}
    \hline
     & $m$& $\psi_{10}$ & $\psi_{11}$ & $\psi_{20}$ & $\psi_{2A}$ & $\psi_{22}$ \\
    \hline
     Sample size 500 &413 & 0.961 & 0.958 & 0.950 & 0.960 & 0.961 \\
     Sample size 2000 &1864 & 0.946 & 0.944 & 0.947 & 0.948 & 0.939 \\
    \hline
    \end{tabular}
\begin{tablenotes}
    \small
    \item Note: $m$, the average resample size in the $m$-out-of-$n$ bootstrap method in 1000 simulated datasets, rounded to integer. For stage-1 parameters ($\psi_{10}, \psi_{11}$), the 95\% confidence intervals \iffalse are\fi were constructed by $m$-out-of-$n$ bootstrap percentiles. For stage-2 parameters ($\psi_{20}, \psi_{2A}, \psi_{22}$), the 95\% confidence intervals \iffalse are\fi were constructed by standard ($n$-out-of-$n$) bootstrap percentiles. 
\end{tablenotes}
  \end{threeparttable}
  \label{WQ-SA_CI_coverage}%
\end{table}%

In summary, Simulation 3 demonstrated that the proposed simulation-based approach could identify a plausible range for the sensitivity parameter. Additionally, it showed that the WQ-SA method was robust and outperformed the naive, CC, and MI methods within a wide range of the sensitivity parameter. Furthermore, it confirmed that the 95\% confidence intervals constructed by the bootstrap percentiles performed well when all the models were correctly specified for the WQ-SA method. 


\subsection{Simulation 4: the impact of an increased number of total stages with nonignorable missing covariates}

To investigate the influence of the nonignorable missing covariates when the number of total stages increases, we designed a simulation study with three stages of intervention. To facilitate comparisons with the 2-stage scenario in Simulation 1, we used the same data generation mechanism for stage 1 and stage 2 as that in Simulation 1, and generated data for stage 3 as follows:  
\begin{align*}
    &X_{3,1} \sim N(0,1); X_{3,2} \sim \text{Uniform}(0,2); \\
    &R_3  \sim \text{Bernoulli} (\{1+\exp(-Y_1-Y_2-0.5X_{3,1}-|X_{3,2}-1|)\}^{-1});\\
    &A_3 \sim 2\times \text{Bernoulli}(\text{expit}(-1+X_{3,1}+X_{3,2})) - 1;\\
    &Y_3 = \beta_{30}+A_3(\psi_{30}+\psi_{3,1}X_{3,2})+\beta_{31}X_{31} +\epsilon_3; \epsilon_3 \sim N(0,1).
\end{align*}
The true values of the parameters $\bm{\theta}_3 = (\psi_{30},\psi_{31},\beta_{30},\beta_{31})$ were $(-1,1,-0.5,0.5)$. $R_3$ was the missingness indicator for the partially observed variable $X_{3,2}$, while $X_{3,1}$ was fully observed. Similar to Simulation 1, we generated 1000 data sets with sample sizes 500 and 2000.  

The final outcome was assumed to be $Y=Y_1+Y_2+Y_3$. The true optimal treatment rule at stage 3 was $2\mathbb{I}(\psi_{30}+\psi_{3,1}X_{3,2} \geq 0)-1$. With the optimal treatment at stage 3, the pseudo-outcome at stage 2 was $Y_{pse,2}=-0.5+|-1+X_{3,2}|+0.5X_{3,1}+Y_1+Y_2$. Because the data generation mechanism at stages 1 and 2 was unchanged from Simulation 1 and $E(Y_{pse,2}\mid \bm{H}_2,A_2,Y_2) = Y_1+Y_2$, the Q-functions, the optimal treatment rules at stages 1 and 2, and $(Y_{pse,1},R_{pse,1})$ were the same as those in Simulation 1. Since $R_3$ only depended on $(X_{3,1},X_{3,2})$, the future-independent missingness assumption held. Among the partially observed variables, only $X_{3,2}$ was employed in the calculation of $Y_{pse,2}$. Therefore, $R_{pse,2} = R_3$. Combining the form of the pseudo-outcomes and the missingness mechanism, we have $P(R_{pse,2}=1\mid \bm{H}_2,A_2,Y_{pse,2}) = \{ 1+\exp(-Y_{pse,2})\}^{-1}$. Since $(X_{2,1},X_{2,2})$ were associated with $Y_{pse,2}$ and independent of $R_{pse,2}$ given $Y_{pse,2}$, we employed $(X_{2,1},X_{2,2})$ as nonresponse instrumental variables when estimating the missingness model for $R_{pse,2}$.  

% \begin{figure}[htbp]
%     \centering
%     \includegraphics[scale = 0.5]{Figures/240427psi10.jpg}
%     \includegraphics[scale = 0.5]{Figures/240427psi11.jpg}\\
%     \includegraphics[scale = 0.5]{Figures/240427psi20.jpg}\includegraphics[scale = 0.5]{Figures/240427psi21.jpg}\\
%     \includegraphics[scale = 0.5]{Figures/240427psi22.jpg}\\
%     \includegraphics[scale = 0.5]{Figures/240427psi30.jpg}
%     \includegraphics[scale = 0.5]{Figures/240427psi31.jpg}\\    
%     \caption{\tb{Boxplots of blip function parameter estimates from the compared methods in Simulation 4. In each boxplot, white and gray boxes are the results for the sample sizes 500 and 2,000, respectively. The horizontal line marks the true values of the blip function parameter estimates. Note that the naive method did not generate estimates of $\psi_{11}$, $\psi_{22}$, and $\psi_{31}$ as $\psi_{11}$, $\psi_{22}$, and $\psi_{31}$ are the coefficients of  partially missing $X_{1,2}$, $X_{2,2}$, and $X_{3,2}$ in the true treatment rules.  }}
%     \label{simu_3stage_fig}
% \end{figure}
% \tp{Not so sure whether to show this figure or not...Our aim is to show that the bias would be larger when the number of stages gets larger, but observing from the bias of $\psi_{10}$ and $\psi_{11}$ estimates, not all of naive, CC, and MI estimates have larger bias than those in Simulation 1.}

% Web Figure \ref{simu_3stage_fig} presents the boxplot of the estimates of blip function parameters $\bm{\psi}_t = (\psi_{10},\psi_{11},\psi_{20},\psi_{2A},\psi_{22},\psi_{30},\psi_{31})$. Since the future-independent missingness held and $(Y_1,Y_2,Y_3)$ were fully observed, the WQ-EE method and the CC method led to unbiased estimates of the 


\begin{table}[htbp]
  \centering
  \caption{Performance of estimated DTRs in Simulation 4}
\begin{threeparttable}
    \begin{tabular}{lcccc}
    \hline
          & \multicolumn{2}{c}{n=500}     & \multicolumn{2}{c}{n=2000} \\
          \cmidrule(lr){2-3}\cmidrule(lr){4-5}
          & Value & Opt\% & Value & Opt\% \\
          \hline
All & 2.983 (0.064) & 0.904 (0.059) & 2.995 (0.031) & 0.950 (0.028) \\
naive & 2.012 (0.064) & 0.255 (0.019) & 2.006 (0.031) & 0.252 (0.009) \\
CC & 2.650 (0.174) & 0.576 (0.114) & 2.622 (0.122) & 0.560 (0.070) \\
MI & 2.348 (0.192) & 0.450 (0.046) & 2.466 (0.103) & 0.479 (0.020) \\
WQ-EE & 2.947 (0.085) & 0.837 (0.097) & 2.983 (0.034) & 0.910 (0.049) \\
    \hline
    \end{tabular}%
\begin{tablenotes}
    \small
    \item Note: Value, the value of the estimated regime; Opt\%, the percentage of patients correctly classified to their true optimal treatments in all the stages;  The numbers outside parentheses represent the sample means, while the numbers inside parentheses represent the sample standard errors.
\end{tablenotes}
  \end{threeparttable}
  \label{simu_3stage_table}%
\end{table}%

Web Table \ref{simu_3stage_table} presents the performance of the methods in comparison in the 3-stage scenario. Compared to Web Table \ref{simutable1} in the 2-stage scenario, the means of the values under the estimated DTRs for all methods decreased. For the All and the WQ-EE method, this decrease was small and it became smaller as the sample size increased. It can be explained by the finite-sample bias accumulated in the backward-induction procedure. In contrast, for the naive, CC, and MI methods, the deterioration of performance was more pronounced, indicating that,  if it is not handled properly, the consequence of the nonignorable missing covariates problem would be more severe when the number of total stages increases. To better understand the reasons behind the performance deterioration, we examined the correct classification rates at different stages. Since we have presented the correct classification rates at stages 1 and 2 in Simulation 1, and the true Q-functions and optimal treatments rules at stages 1 and 2 in Simulation 1 and Simulation 4 were the same, we examined the correct classification rates at stage 3 and at stages 1 and 2, respectively. 


\begin{table}[htbp]
  \centering
  \caption{Correct classification rates in Simulation 4}
  \begin{threeparttable}
    \begin{tabular}{lcccc}
    \hline
          & \multicolumn{2}{c}{n=500}     & \multicolumn{2}{c}{n=2000} \\
          \cmidrule(lr){2-3}\cmidrule(lr){4-5}
          & Stages 1\&2 & Stage 3 & Stages 1\&2 & Stage 3 \\
          \hline
All & 0.919  (0.060) & 0.980  (0.017) & 0.958  (0.027) & 0.990  (0.008) \\
naive & 0.501  (0.021) & 0.506  (0.021) & 0.501  (0.011) & 0.504  (0.010) \\
CC & 0.594  (0.117) & 0.958  (0.035) & 0.568  (0.067) & 0.981  (0.016) \\
MI & 0.485  (0.040) & 0.931  (0.058) & 0.506  (0.016) & 0.936  (0.026) \\
WQ-EE &0.873  (0.098) & 0.958  (0.035) & 0.928  (0.054) & 0.981  (0.016) \\
    \hline
    \end{tabular}%
\begin{tablenotes}
    \small
    \item Note: Stages 1\&2, the correct classification rate at stages 1 and 2; Stage 3, the correct classification rate at stage 3; The numbers outside parentheses represent the sample means, while the numbers inside parentheses represent the sample standard errors.
\end{tablenotes}
  \end{threeparttable}
  \label{simu_3stage_table2}%
\end{table}%

We used the results in  Web Table \ref{simu_3stage_table2} to investigate why    the increased number of total stages  led to  lower overall correct classification rate if nonignorable missing covariate were  not handled properly. For the naive method, its correct classification rate at stages 1 and 2 in Simulation 4 was almost the same as that in Simulation 1 (Web Table~\ref{simutable1}), and the decrease in its overall correct classification rate in Web Table~\ref{simu_3stage_table} was mainly explained by its poor average correct classification rate  at stage 3, which was near $0.5$. In contrast,  for the CC and the MI method, the average correct classification rates at stage 3 was higher than $0.93$, but their average correct classification rates at stages 1 and 2 dropped a lot compared to those in Simulation 1  (Web Table \ref{simutable1}). And this decrease was due to the bias from nonignorable missing pseudo-outcomes accumulated over more stages in the backward induction procedure.     

% \comment{Maybe a couple of sentences to summarize Simulation 4}

% \tb{Generally, the increased number of total stages would lead to lower overall correct classification rates when covariates were nonignorable missing. If the nonignorable missingness was properly addressed, the decrease of overall correct classification rate would be low and can be mitigated by increasing sample size. However, if the nonignorable missingness was not properly handled, the decrease in overall correct classification would obvious since the bias would accumulate in the backward-induction procedure.  }

In summary, an increased number of total stages would generally lead to lower overall correct classification rates of the estimated optimal DTRs when covariates had nonignorable missing values. If the nonignorable missingness was properly addressed, the decrease in the overall correct classification rates would be small and could be mitigated by increasing the sample size. However, if the nonignorable missingness was not properly handled, the decrease in overall correct classification rates could be substantial due to the bias accumulated in the backward-induction procedure.




\section{Web Appendix I: Additional details  of  the MIMIC-III data and the estimation procedure}


The proposed methods were applied to the electronic medical records and administrative data from the MIMIC-III database \citep{Johnson_et_al_2016}. MIMIC-III is a retrospectively collected and freely available database accessible through PhysioNet \citep{Goldberger_et_al_2000}. The MIMIC-III database contains de-identified and anonymized health information from over 45,000 patients who were treated in the Intensive Care Units (ICUs) of Beth Israel Deaconess Medical Center in Boston, Massachusetts. We utilized similar patient selection criteria as in \cite{Speth_et_al_2022}. The detailed patient selection criteria are presented in Web Figure~\ref{patient_select}.

\begin{figure}[htbp]
    \centering
    \includegraphics[scale=0.75]{patient_select2.png}
    \caption{Eligibility criteria and patient population for the selected cohort from the MIMIC-III database.}
    \label{patient_select}
\end{figure}

As depicted in Web Figure \ref{patient_select}, This study included adult patients (age $\geq$ 18) diagnosed with acute emergent sepsis who were transferred from the Emergency Department (ED) and admitted to the medical ICU. Patients must be sent to the medical ICU within 1 hour after hospital admission and had documented fluid intake within the first hour after hospital admission. Furthermore, following \cite{Speth_et_al_2022}, to ensure a relatively homogeneous patient cohort, we only included patients who survived for a minimum of 48 hours after ICU admission. 

Since there were very few patients with hemodynamic variables recorded precisely at 180 minutes after ICU admission, we calculated the mean of the hemodynamic variables recorded within a half-hour window around the third hour. That is, we used the data collected between 2.5 and 3.5 hours after ICU admission as surrogates for the values of the hemodynamic variables three hours after ICU admission. If a hemodynamic variable had no record during this specific time interval, we considered the corresponding record three hours after admission to be missing. Under this criterion, the missing data proportion of hemodynamic variables is provided in Web Table~\ref{missingtable}.

\begin{table}[htbp]
    \centering
    \caption{The missing proportion of hemodynamic variables 3 hours after ICU admission in the MIMIC-III data}
    \begin{tabular}{lc}
    \hline
        covariates &  missing proportion\\
        \hline 
        heartrate at 3 hours & 7.50 \% \\
        mean blood pressure at 3 hours & 10.2 \%  \\
        SpO2 at 3 hours  & 10.4 \% \\
        respiratory rate at 3 hours & 10.2 \% \\
        temperature at 3 hours & 19.8 \% \\
        urine output in 0-3 hours & 0.00 \% \\
    \hline
    \end{tabular}
    \label{missingtable}
\end{table}

%To assess the out-of-sample performance of the estimated optimal DTRs, we employed a 5-fold cross-validation. For each bootstrap sample, the available dataset was divided into 5 folds. 
%To estimate the mean of the out-of-sample improvement of SOFA at 24h, we employed a 5-fold cross-validation \tb{within ordinary $n$-out-of-$n$ bootstrap}. For each bootstrap sample, the available dataset was divided into 5 folds. Four folds were combined to create a training dataset, while the remaining 5th fold was used as the testing dataset. We alternated the roles of the folds so that each fold served as the test set exactly once. For each fold, we re-estimated the semiparametric missingness probability model and the Q-functions based on the training set. The expected outcome on the testing set, under the estimated DTR, was then computed using the estimated missingness probability models and Q-functions. It should be noted that the optimal DTR may differ across different folds. 

Following \cite{xu2015regularized}, we performed cross-validation and $m$-out-of-$n$ bootstrap to assess the performance of the estimated optimal DTRs. Specifically, to estimate the expected out-of-sample improvement of SOFA score at 24 hours post-admission, we randomly split the MIMIC-III data for 1000 times. For each split, we used $4/5$ of the data (rounded to the nearest integer) as the training set and the remaining data as the testing set. 
The semiparametric missingness probability model and the Q-functions were re-estimated based on the training set of that split. The expected outcome in the testing set, under the estimated DTR, was then computed using the estimated missingness probability models and Q-functions. It should be noted that the optimal DTR may differ across different splits. Nevertheless, the average of the expected optimal outcomes in the testing sets from repeated splitting can be interpreted as the expected SOFA score at 24 hours after ICU admission that would have been obtained with out-of-sample data under the estimated regimes.

For the inference of the improvement of SOFA score at 24 hours post-admission, we further calculated the bootstrap standard deviation and the 95\% confidence interval  %$\widehat{E}[Q_1\{\bm{h}_1,\widehat{d}(\bm{h}_1); \widehat{\bm{\theta}}_1\} ] - \widehat{E}(Y)$ 
using the $m$-out-of-$n$ bootstrap method based on the original data (i.e., without cross-validation). We used the double bootstrap method described in Web Appendix G to determine the bootstrap resample size $m$, with the number of bootstrap replications set as $B_1=B_2=200$. Specifically, when selecting $m$ and conducting $m$-out-of-$n$ bootstrap for the MI method, we followed the suggestion from \cite{Shen_Hubbard_Linn_2023} and used 25 imputations.  

%\baselineskip=20pt
\bibliographystyle{biom} 
\bibliography{DTR_MNAR_bib}



\end{document}



